\documentclass[11pt,letterpaper]{article}
\usepackage[T1]{fontenc}
\usepackage[utf8]{inputenc}
\usepackage{lmodern}
\usepackage[margin=1in]{geometry}
\usepackage{amsmath,amssymb,amsthm,mathtools}
\usepackage{aliascnt}
\usepackage{array,booktabs,longtable}
\usepackage{enumitem,microtype,xcolor}
\usepackage{xurl}
\usepackage[colorlinks=true,linkcolor=blue!45!black,citecolor=blue!45!black,urlcolor=blue!45!black]{hyperref}
\usepackage[nameinlink,noabbrev]{cleveref}
\usepackage{listings}
\usepackage{fancyhdr}
\pagestyle{fancy}
\fancyhf{}
\fancyhead[L]{\small Positive factorial-window determinants}
\fancyhead[R]{\small Journal version}
\fancyfoot[C]{\thepage}
\setlength{\headheight}{14pt}
\renewcommand{\sectionmark}[1]{\markboth{\thesection.\ #1}{}}
\setlength{\parskip}{.35em}
\setlength{\parindent}{1.1em}
\setlength{\emergencystretch}{3em}
\setlist{topsep=4pt,itemsep=3pt,parsep=1pt}
\newtheorem{theorem}{Theorem}[section]
\newaliascnt{lemma}{theorem}
\newtheorem{lemma}[lemma]{Lemma}
\aliascntresetthe{lemma}
\newaliascnt{proposition}{theorem}
\newtheorem{proposition}[proposition]{Proposition}
\aliascntresetthe{proposition}
\newaliascnt{corollary}{theorem}
\newtheorem{corollary}[corollary]{Corollary}
\aliascntresetthe{corollary}
\theoremstyle{definition}
\newaliascnt{definition}{theorem}
\newtheorem{definition}[definition]{Definition}
\aliascntresetthe{definition}
\newaliascnt{remark}{theorem}
\newtheorem{remark}[remark]{Remark}
\aliascntresetthe{remark}
\newtheorem*{example}{Example}
\crefname{lemma}{lemma}{lemmas}
\Crefname{lemma}{Lemma}{Lemmas}
\crefname{proposition}{proposition}{propositions}
\Crefname{proposition}{Proposition}{Propositions}
\crefname{corollary}{corollary}{corollaries}
\Crefname{corollary}{Corollary}{Corollaries}
\crefname{definition}{definition}{definitions}
\Crefname{definition}{Definition}{Definitions}
\crefname{remark}{remark}{remarks}
\Crefname{remark}{Remark}{Remarks}
\newcommand{\Q}{\mathbb Q}
\newcommand{\Z}{\mathbb Z}
\newcommand{\R}{\mathbb R}
\newcommand{\C}{\mathbb C}
\newcommand{\Fp}{\mathbb F_p}
\newcommand{\Zp}{\mathbb Z_p}
\newcommand{\Qp}{\mathbb Q_p}
\newcommand{\dd}{\,\mathrm d}
\newcommand{\vp}{v_p}
\newcommand{\J}{\mathcal J}
\newcommand{\Bfun}{\mathcal B}
\newcommand{\arith}{\mathrm{arith}}
\newcommand{\an}{\mathrm{an}}
\newcommand{\pos}[1]{\left(#1\right)_{+}}
\newcommand{\floor}[1]{\left\lfloor #1\right\rfloor}
\newcommand{\ceil}[1]{\left\lceil #1\right\rceil}
\newcommand{\abs}[1]{\left|#1\right|}
\newcommand{\norm}[1]{\left\|#1\right\|}
\DeclareMathOperator{\rank}{rank}
\DeclareMathOperator{\spanQ}{span}
\DeclareMathOperator{\arcosh}{arcosh}
\DeclareMathOperator{\supp}{supp}
\lstset{language=Python,basicstyle=\ttfamily\scriptsize,columns=fullflexible,
  keepspaces=true,breaklines=true,breakatwhitespace=false,showstringspaces=false,
  frame=single,framerule=.3pt,rulecolor=\color{black!25},
  numbers=left,numberstyle=\tiny\color{black!50},numbersep=6pt,
  xleftmargin=1.5em,framexleftmargin=1.3em,tabsize=4}
\hypersetup{pdfauthor={Zeta Irrationality Project},pdfsubject={A complete candidate proof and its mathematical background}}

\title{\textbf{Positive factorial-window determinants
and the irrationality of $\zeta(5)$}}
\author{Zeta Irrationality Project}
\date{13 November 2026}
\hypersetup{pdftitle={Positive factorial-window determinants and the irrationality of zeta(5)}}
\begin{document}
\maketitle
\begin{center}
\small Draft for mathematical review. The exact rational comparison has been
reproduced; this argument has not been formally verified in a proof assistant.
\end{center}
\begin{abstract}
We present an internally audited candidate proof of the irrationality of
$\zeta(5)$ using a positive Gram determinant indexed by a consecutive window
of factorial-normalized polynomials. Rational moments and an affine
Stieltjes formula produce a polynomial in $\zeta(5)$ of degree at most $N$.
A local polynomial-restriction argument and a common-root improvement bound
its primitive coefficient normalizer. An explicit positive logarithmic
potential, together with quantified polynomial discretization, bounds the
determinant itself. For the window $m=N/20$, $R=5N/4$, the arithmetic
logarithmic rate is below $5.4101$ and the analytic decay rate exceeds
$5.4133$, on the scale $N^2$. A complete integer-and-rational certificate
verifies the strict gap. The historical introduction records classical and
recent results through 13 November 2026, including public preprints and
formalization repositories, with their verified software versions and stated
trust boundaries. A preceding $\zeta(3)$ specialization gives an unshifted-window warm-up with exact positive gap $(9/2)\log 2-3$. A final section derives the higher-odd-value functional, quantifies the failure of the present bounds at $\zeta(7)$, and identifies the improvements needed beyond it.

\end{abstract}
\clearpage
\begingroup
\footnotesize
\setlength{\parskip}{0pt}
\tableofcontents
\endgroup
\clearpage

% BEGIN SHARED CORE 01_history
\section{Historical introduction}\label{sec:history}

For a real number $s>1$, the Riemann zeta function is
\[
  \zeta(s)=\sum_{n=1}^{\infty}\frac{1}{n^s}.
\]
This introduction concerns its values at ordinary positive integers; it does
not identify them with the distinct objects called $p$-adic zeta values.
A real number is \emph{irrational} if it does not belong to $\Q$, and
\emph{transcendental} if it is not a root of a nonzero polynomial with rational
coefficients.  Euler's evaluation of $\zeta(2)$, and more generally of the
positive even values as nonzero rational multiples of powers of $\pi$, gave
the subject its first decisive examples.  Euler's paper E41 was written in
1735 and published in 1740; Lindemann's 1882 transcendence theorem for $\pi$
subsequently made every positive even zeta value transcendental
\cite{euler1740,lindemann1882,dlmf}.

The odd values required a different approach.  Ap\'ery announced the
irrationality of $\zeta(3)$ in 1978 and published his proof in 1979.
Beukers' integral proof appeared in the same year
\cite{apery1979,beukers1979}.  Later advances developed in two directions:
proving that many odd zeta values are irrational, without specifying which
ones, and shrinking explicit finite lists that cannot consist entirely of
rational numbers.  Neither kind of statement automatically identifies an
individual irrational value.

Here is notation that makes these distinctions precise.  For odd integers
$3\leq a\leq b$, put
\[
 W(a,b)=\{\zeta(a),\zeta(a+2),\ldots,\zeta(b)\}.
\]
Saying that $W(a,b)$ contains an irrational number means that at least one
of these numbers is irrational; it does not assert that all of them are.
For an integer $s\geq3$, define
\begin{align*}
 N(s)&=\#\{k:3\leq k\leq s,\ k\text{ odd},\ \zeta(k)\notin\Q\},\\
 D(s)&=\dim_{\Q}\operatorname{span}_{\Q}
       \bigl(\{1\}\cup\{\zeta(k):3\leq k\leq s,\ k\text{ odd}\}\bigr).
\end{align*}
The symbol $\#$ denotes the number of elements of a finite set.  A family
$x_1,\ldots,x_m$ is \emph{linearly independent over $\Q$} when
$\sum_i q_i x_i=0$ with $q_i\in\Q$ forces every $q_i=0$; its rational span
is the set of its rational linear combinations.  Its dimension counts a
maximal independent subfamily.  In particular,
\[
                    D(s)\leq N(s)+1.
\]
Thus a dimension lower bound implies a counting lower bound after
subtracting one, whereas many irrational values can still have a small
rational span.  All logarithms below are natural logarithms.  Bounds
containing $\varepsilon$ hold for every fixed positive $\varepsilon$ and
all sufficiently large admissible $s$.

\subsection{Mathematical milestones through 13 November 2026}

\Cref{tab:history-results} records the principal milestones relevant to
these two directions, together with recent public developments that have
not yet acquired the same publication status as the classical results.
Dates distinguish a preprint's appearance from its journal publication
where this matters.  A repository is cited as a repository: its presence
in this table is not a claim of journal peer review.  The formal status
of the recent developments is specified separately in
\cref{tab:history-formal}.

\begingroup
\footnotesize
\renewcommand{\arraystretch}{1.18}
\setlength{\LTpre}{6pt}
\setlength{\LTpost}{6pt}
\begin{longtable}{@{}>{\raggedright\arraybackslash}p{.14\textwidth}>{\raggedright\arraybackslash}p{.57\textwidth}>{\raggedright\arraybackslash}p{.21\textwidth}@{}}
\caption{Selected mathematical milestones.  ``One in $W$'' means at least
one irrational member; ``two in $W$'' means at least two.  Recent source
claims and conditional developments are explicitly identified.}
\label{tab:history-results}\\
\toprule
Date & Mathematical statement & Source and status\\
\midrule
\endfirsthead
\toprule
Date & Mathematical statement & Source and status\\
\midrule
\endhead
\midrule
\multicolumn{3}{r}{\emph{Continued on the next page}}\\
\endfoot
\bottomrule
\endlastfoot
1735/1740; 1882 & Euler evaluates the positive even values; Lindemann's
theorem makes every $\zeta(2m)$, $m\geq1$, transcendental. &
\cite{euler1740,lindemann1882,dlmf}.\\
1978/1979 & $\zeta(3)$ is irrational.  Beukers supplies a proof using
multiple integrals. & Ap\'ery and Beukers; published
\cite{apery1979,beukers1979}.\\
2000/2001 & Infinitely many odd values are irrational; in particular
$D(s)\geq (1-\varepsilon)\log s/(1+\log2)$ for large odd $s$. &
Rivoal; Ball--Rivoal; published \cite{rivoal2000,ballrivoal2001}.\\
2001/2002 & One in $W(5,21)$ (nine values). & Rivoal's 2001 preprint,
2002 journal article \cite{rivoal2002}.\\
2001 & One in $W(5,11)$ (four values).  A separate paper gives one in
$W(7,35)$ and one in $W(9,51)$. & Zudilin; published
\cite{zudilin2001,zudilin2001windows}.\\
2018/2019 & $N(s)\geq 2^{(1-\varepsilon)\log s/\log\log s}$. &
Fischler--Sprang--Zudilin; published \cite{fsz2019}.\\
2018/2020 & Two in $W(5,69)$ (33 values). & Rivoal--Zudilin;
published \cite{rivoalzudilin2020}.\\
2019/2020 & $N(s)\geq(1.192507\ldots-\varepsilon)\sqrt{s/\log s}$. &
Lai--Yu; published \cite{laiyu2020}.\\
2021/2022 & Two in $W(5,35)$ (16 values). & Lai--Zhou;
published \cite{laizhou2022}.\\
20 Apr. 2026 & $D(s)\geq0.21\sqrt{s/\log s}$ for sufficiently large odd
$s$.  This is a dimension bound, not merely a count. & Fischler;
published \cite{fischler2026}.\\
16 Sept. 2026 & $N(s-1)\geq1.284\sqrt{s/\log s}$ for sufficiently large
even $s$. & Li Lai; 2025 preprint, 2026 journal publication
\cite{lai2026count}.\\
17 Sept. 2026 & $\zeta(5)$ is irrational; the preprint also states an
irrationality-measure estimate.  The implementations below certify the
irrationality endpoint, not that additional estimate. & Aabir Fauzan;
Zenodo v1 \cite{fauzan2026}.\\
Sept.--Oct. 2026 & One in $W(7,21)$ (eight values), in separate public
developments. & Devincenzis, Lin, and Krisman; repositories/preprint
\cite{devincenzis2026,lin2026,krisman2026}.\\
29 Sept. 2026 & One in $W(9,33)$ (13 values).  Also, for every odd
$r\geq3$, some odd-indexed $\zeta(k)$ with $r+2\leq k\leq6r-1$ is
irrational. & Yongxi Lin; repository \cite{lin2026}.\\
29--30 Sept. 2026 & $1,\zeta(2),\zeta(3)$ are linearly independent over
$\Q$. & Qian Tang; repository \cite{tang2026}.\\
6 Oct. 2026 & One in $W(7,19)$ (seven values). & Jingwei Liu;
Zenodo/repository v1.0.0 \cite{jingweiliu2026}.\\
8 Oct. 2026 & $D(s-1)\geq1.119\log s/(1+\log2)$ for sufficiently large
even $s$.  The paper also proves a distinct $p$-adic analogue. & Li Lai;
published; real bound superseded by Fischler's growth rate
\cite{lai2026small}.\\
5/15 Oct. 2026 & $1,\zeta(3),\zeta(4)$ are linearly independent over
$\Q$; hence $\zeta(3)/\pi^4$ is irrational. & Li Lai;
arXiv v2, no formal verification identified \cite{lai2026linear}.\\
15--16 Oct. 2026 & Proposed four-value conclusion: one in $W(7,13)$.
The public Lean endpoint depends on eight numerical admissions; the
certificate-conditional theorem is separately proved. & Morley Davidson;
public project snapshot \cite{davidson2026}.\\
Nov. 2026 & A second verification of the four-value conclusion, including
the finite numerical checks, is reported.  Its source was unavailable for
inspection. & Zaz Brown; unpublished/private development reported by the
project author; see the qualification below.\\
\end{longtable}
\endgroup

The shifted-window numbers in the 2001 row refer to Zudilin's published
\emph{Mathematical Notes} paper.  Its bounds $35$ and $51$ improve the
$37$ and $53$ in the earlier arXiv note \cite{zudilin2001note}.
These are different versions of the historical record and should not be
silently interchanged.  A separate September 2026 preprint by Prince
Anand claims the irrationality of $\zeta(7)$; a public issue identifies a
specific arithmetic objection to its argument.  We record that claim
and objection without treating the claimed theorem as established
\cite{anand2026,anandobjection}.  The table is a sourced account of major
results and the located recent records, not an assertion that every
unpublished announcement has been discovered.

\subsection{Formal verification: statements, versions, and assumptions}

A formal proof is an expression checked against explicit inference rules
by a proof assistant's \emph{kernel}.  The \emph{endpoint} is the precise
proposition that the development ultimately proves.  Inspecting it is
essential: a proof of an implication from numerical hypotheses has a
different scope from a proof that also establishes those hypotheses.
An \emph{axiom audit} lists the assumptions on which that endpoint depends.
In Lean, \texttt{sorryAx} is the assumption introduced by an admitted proof
step; the presence of an unused admission elsewhere in a repository does
not by itself put that admission in the endpoint's dependency chain.

In \cref{tab:history-formal}, ``standard Lean axioms'' abbreviates
\texttt{propext}, \texttt{Classical.choice}, and \texttt{Quot.sound}:
propositional extensionality, classical choice, and quotient soundness.
It does not mean an absence of foundational assumptions.  An additional
prime-number-theorem axiom is a mathematical assumption beyond this
standard foundation.  An \emph{oracle} is a trusted procedure allowed to
introduce a theorem without an ordinary proof term for each computation.
An independent checker can provide useful additional assurance, but its
acceptance still concerns the submitted endpoint and assumptions.

The inventory below is based on primary papers, source files, declared
toolchains, and the build or audit records explicitly identified in the
references.  \emph{No development was rebuilt or replayed for this
historical survey.}  ``Reports'' therefore attributes a verification
claim to its published record.  A commit identifier fixes a source
snapshot; a software version fixes the named implementation release.
Unknown versions are marked unknown instead of being reconstructed from
an unrelated current website header.

\begingroup
\footnotesize
\renewcommand{\arraystretch}{1.18}
\begin{longtable}{@{}>{\raggedright\arraybackslash}p{.24\textwidth}>{\raggedright\arraybackslash}p{.22\textwidth}>{\raggedright\arraybackslash}p{.46\textwidth}@{}}
\caption{Located formal implementations and their documented scope.
Source snapshots and dependency pins are supplied in the references.
This table distinguishes completed author-reported checks, explicit
assumptions, admissions, and unavailable evidence.}
\label{tab:history-formal}\\
\toprule
Development & Proof software/version & Endpoint and documented qualification\\
\midrule
\endfirsthead
\toprule
Development & Proof software/version & Endpoint and documented qualification\\
\midrule
\endhead
\midrule
\multicolumn{3}{r}{\emph{Continued on the next page}}\\
\endfoot
\bottomrule
\endlastfoot
Chyzak--Mahboubi--Sibut-Pinote--Tassi (2014); Mahboubi--Sibut-Pinote (2021)
\cite{chyzak2014,mahboubi2021,aperysource} & Coq; completed paper cites
8.11.0.  Current source requires Coq 8.20+; CI also configures Rocq 9.0/9.1.
& $\zeta(3)$ irrational.  The 2014 account left a weak prime-number-theorem
consequence open; the 2021 account reports the complete constructive
proof.  Computer-algebra guesses are subsequently checked in Coq.\\
Manuel Eberl, AFP, 27 Dec. 2019 \cite{eberl2019} & Isabelle/HOL;
original release unverified here.  Inspected development build:
Isabelle revision \texttt{559e5864efb7}. & $\zeta(3)$ irrational, following
Beukers.  The inspected successful AFP build is dated 13 Nov. 2026, with
AFP revision \texttt{62317ad25318}; these are not claimed as the 2019
versions.\\
Junqi Liu--Jujian Zhang--Lihong Zhi \cite{liuzhangzhi2025} & Lean 4.18.0;
Mathlib v4.18.0. & Authors report a complete proof of $\zeta(3)$
irrational.  Endpoint source and dependency pins were inspected;
no new transitive axiom audit was run.\\
C\'esar Del Solar \cite{delsolar2026} & Lean 4.35.0-rc2. & Fauzan's
$\zeta(5)$ irrationality.  Reports only standard Lean axioms for
\texttt{zeta5\_is\_irrational}; four unused admissions lie outside its
dependency chain.  The prime number theorem is imported as a proof.\\
Moritz Firsching \cite{firsching2026} & Lean 4.34.0-rc1. & Fauzan's
$\zeta(5)$ irrationality.  Reports only standard Lean axioms for
\texttt{irrational\_five}; unrelated imported admissions do not occur in
that dependency chain.  A successful build for the pinned snapshot is
recorded by Reservoir.\\
Dan Romik \cite{romik2026} & Lean 4.34.0; Mathlib v4.34.0. & Fauzan's
$\zeta(5)$ irrationality \emph{assuming an explicit prime-number-theorem
axiom}.  Reports no admissions, but the mathematical axiom
\texttt{chebyshev\_theta\_asymptotic} remains.\\
Alex Raymond Hyman, AFP, 26 Oct. 2026 \cite{hyman2026} & Isabelle/HOL;
exact release not established from the inspected entry. & Port of
Del Solar's proof of $\zeta(5)$ irrationality.  Prime number theorem from
AFP; the \texttt{holds} oracle is used for 72 numerical inequalities.
This is a port, not an independent mathematical proof.\\
Gennaro Marco Devincenzis \cite{devincenzis2026} & Lean 4.35.0-rc2. &
One in $W(7,21)$.  Reports standard Lean axioms and kernel-checked finite
calculations.  The repository reports Lean, NanoDa, and \texttt{con-ron}
checks using the tools bundled with that Lean release.\\
Yongxi Lin \cite{lin2026} & Lean 4.35.0-rc3. & Eight-value, 13-value,
and uniform-window endpoints in the preceding table.  Reports standard
Lean axioms.  Lean, NanoDa, and \texttt{con-ron} are the versions bundled
with this release; no separate checker version is declared.\\
Qian Tang \cite{tang2026} & Lean 4.35.0-rc3. & Rational linear
independence of $1,\zeta(2),\zeta(3)$.  Reports standard Lean axioms and
checks with Lean, NanoDa, and \texttt{con-ron}; independent release
identifiers for the latter checkers were not established.\\
Jingwei Liu \cite{jingweiliu2026} & Lean 4.34.0-rc2. & One in
$W(7,19)$.  Reports a fresh project rebuild, standard Lean axioms, and
kernel verification of the finite numerical work.  Jingwei Liu is not
the Junqi Liu of the earlier Ap\'ery formalization.\\
Shawn Krisman \cite{krisman2026} & Lean 4.35.0; Mathlib v4.35.0. &
One in $W(7,21)$.  Reports standard Lean axioms and kernel-checked
numerical certificates; completed 22 Oct. 2026.  The repository does
not claim expert human review.\\
Morley Davidson \cite{davidson2026} & Lean 4.33.1.  External numerical
checks: python-flint 0.9.0, Arb at 256-bit precision. & Four-value
certificate-conditional theorem reports standard Lean axioms.  Its
supplied-certificate wrapper still depends on eight numerical
admissions and hence on \texttt{sorryAx}.  External interval checks
are not Lean proof terms.\\
Zaz Brown (reported November 2026) & Lean; version, source commit,
and dependency versions unavailable. & The project author reports a
second four-value verification on GitLab that includes all finite
checks.  No accessible source, endpoint, axiom report, or build receipt
was obtained.  This row records the report, not an independently
inspected certification.\\
\end{longtable}
\endgroup

There are therefore three located Lean implementations of Fauzan's
argument, with different documented assumptions, and an additional
Isabelle/HOL port.  Romik's explicit prime-number-theorem assumption
must remain visible when counting these implementations as
certifications.  Likewise, the public four-value endpoint
\path{OddZeta.FourValue.irrational_zeta7_or_zeta9_or_zeta11_or_zeta13}
must be distinguished from its fully proved implication from
certificate hypotheses \cite{davidson2026}.

The report concerning Zaz Brown was communicated by the project author
on 13 November 2026, with the qualification that the GitLab development
might not yet be public.  We preserve that report without inventing a
URL, a release number, or a publication date.  Its asserted coverage of
the finite checks cannot yet replace an inspected proof object in this
inventory.  These distinctions concern the evidence available for each
formal development; they do not supply a proof of any new mathematical
step in the present manuscript.

\subsection{Methodological precedents}

Determinants already have a substantial place in irrationality research.
Zudilin's determinantal criterion extends the classical argument from
small nonzero integer linear forms by exploiting additional structure
\cite{zudilin2017}.  Haynes and Zudilin study determinants of matrices
whose entries are zeta values along arithmetic progressions and derive
Diophantine consequences from their asymptotics \cite{hayneszudilin2015}.
A \emph{Hankel matrix} has entries of the form $m_{i+j}$ for a sequence
$(m_j)$; its entries are constant along each anti-diagonal.  Fauzan's
2026 construction also uses Hankel determinants \cite{fauzan2026}.
Accordingly, a self-contained derivation of a determinant argument is
not, by itself, a claim that the method is historically new.  The role
of the sections that follow is to specify the proposed construction,
its exact hypotheses, and the inequalities needed to complete it.

% END SHARED CORE 01_history

% BEGIN JOURNAL ADDITION zeta3
\section{A determinantal warm-up: the irrationality of \texorpdfstring{$\zeta(3)$}{zeta(3)}}
\label{sec:zeta3-warmup}

The construction used below for $\zeta(5)$ has a simpler specialization to
$\zeta(3)$. We give it first to expose the common mechanism: a positive
Gram determinant is a rational polynomial in the target number; an
arithmetic estimate clears its coefficients; and an analytic estimate
forces its value, after clearing, to approach zero. In this specialization
the comparison of the two estimates reduces to $\log2>2/3$. No finite
numerical certificate is needed.

\subsection{Relation to the classical proofs}

Ap\'ery's proof constructs rational linear forms in $\zeta(3)$ and combines
their arithmetic properties with sufficiently rapid decay
\cite{apery1979}. Beukers supplies a positive integral realization of
suitable forms, making their nonvanishing and their smallness accessible
through the integrand \cite{beukers1979}. The argument here retains this
arithmetic--analytic division of labor. Its positive object is a Gram
determinant, and its polynomial degree grows with the matrix size. Both
its denominator cost and its analytic saving are estimated on the scale
$N^2$; substituting a hypothetical rational value with denominator $b$
costs only $N\log b$.

There are substantial precedents for this formulation. Pr\'evost gives
Pad\'e-based proofs for $\zeta(2)$ and $\zeta(3)$
\cite{prevost1996}; his later positive Stieltjes representation and the
subsequent work with Rivoal develop the associated rational approximation
framework \cite{prevost2019,prevostrivoal2021}. Zudilin's determinantal
criterion explains how positive moment determinants can strengthen an
irrationality argument \cite{zudilin2017}. The positive-functional setting
also appears in Calegari's exposition, which explicitly treats determinants
with several simple poles and includes the present rational-weight family
\cite{calegari2026}. Accordingly,
this section is an explicit specialization and comparison, with no claim
of novelty for the general method or of a new formal verification.

The connection to Pr\'evost is more specific than the use of positivity.
His construction uses rational Bernoulli moments and square-node products
of the form
\[
 \binom{x-1}{k}\binom{x+k}{k}
 =\frac{\prod_{a=1}^{k}(x^2-a^2)}{(k!)^2}.
\]
With $x=i\sqrt t$, this becomes $(-1)^kD_k(t)/(k!)^2$, where $D_k$
is defined below. A particular subfamily, denoted
$RPA(m,2m-1,2m)$ in the arXiv version of \cite[Section~5.2]{prevost2019},
recovers Ap\'ery's numbers. That identification applies to the stated
indices; it does not identify every determinant built from this moment
functional with the classical linear forms.

\begin{theorem}\label{thm:z3-warmup}
The number $\zeta(3)=\sum_{n\ge1}n^{-3}$ is irrational.
\end{theorem}

All integers $N\ge1$ are admissible in this section. Our choice of window
is $m=0$, $R=N$, so there is no omitted initial block and no common-root
improvement to estimate. The subscript $3$ in $\mu_3$ and
$\mathcal J_{3,X}$ denotes the target zeta argument and distinguishes
these objects from the subsequent $\zeta(5)$ construction. The elementary
interpolation and potential calculations used here are proved in detail
in Sections~\ref{sec:integer-valued}--\ref{sec:local-kernel} and
\ref{sec:potential}--\ref{sec:discretization}; we record explicitly the
changes needed at the present parameters.

\subsection{A positive functional containing only \texorpdfstring{$\zeta(3)$}{zeta(3)}}

Put $f(x)=(e^{2\pi x}-1)^{-1}$ for $x>0$. Define a measure $\mu_3$ in
the variable $t=x^2$ by
\begin{equation}\label{eq:z3-measure}
 \int g(t)\,\dd\mu_3(t)
 =\int_0^\infty g(x^2)x^3 f''(x)\,\dd x
\end{equation}
for nonnegative measurable $g$. The density is strictly positive, since
$f''(x)=(2\pi)^2\sum_{n\ge1}n^2e^{-2\pi nx}$. Its value at zero,
defined by continuity, is $1/\pi$, and
\begin{equation}\label{eq:z3-density-bound}
 0<x^3f''(x)\le C(1+x^3)e^{-2\pi x}\qquad(x>0).
\end{equation}
The Laurent expansion of $f$ proves the assertions at zero; the positive
exponential series proves the bound for $x\ge1$.

Use the Bernoulli convention
$z/(e^z-1)=\sum_{n\ge0}\mathrm B_nz^n/n!$. Integration of the positive
exponential series and Euler's even-zeta formula give
\begin{equation}\label{eq:z3-moments}
 \int t^j\,\dd\mu_3(t)
 =\frac{(2j+3)!}{(2\pi)^{2j+2}}\zeta(2j+2)
 =\frac{(-1)^j(2j+3)}2\mathrm B_{2j+2}\in\Q.
\end{equation}
In particular, $\int\dd\mu_3=1/4$.
For $y>0$ the corresponding Stieltjes transform is
\begin{equation}\label{eq:z3-stieltjes}
 \int\frac{\dd\mu_3(t)}{t+y^2}
 =y^2\zeta(3,y)-\frac1{2y}-\frac12,
 \qquad \zeta(3,y)=\sum_{n\ge0}(n+y)^{-3}.
\end{equation}
Here is a direct derivation. If
$K(y)=\int_0^\infty xf(x)/(x^2+y^2)\,\dd x$, Binet's identity
\eqref{eq:binet-digamma} reads
$2K(y)=\log y-(2y)^{-1}-\psi(y)$.
Two integrations by parts in the left side of
\eqref{eq:z3-stieltjes}, using
$x^3/(x^2+y^2)=x-y^2x/(x^2+y^2)$ and the harmonicity of
$x/(x^2+y^2)$, give $y^2K''(y)$. Differentiation of Binet's identity
gives $K''(y)=\zeta(3,y)-(2y^3)^{-1}-(2y^2)^{-1}$.
The integrations by parts have vanishing boundary terms: at zero each
boundary product is $O(x)$, and at infinity the derivatives of $f$
decay exponentially. This proves \eqref{eq:z3-stieltjes}.

Write $H_n^{(3)}=\sum_{i=1}^n i^{-3}$, with $H_0^{(3)}=0$. On rational
functions having only simple poles among $-1,-4,-9,\ldots$, define the
$\Q$-linear functional $\mathcal J_{3,X}$ by
\begin{align}
 \mathcal J_{3,X}(t^j)
   &=\frac{(-1)^j(2j+3)}2\mathrm B_{2j+2},\notag\\
 \mathcal J_{3,X}\left(\frac1{t+i^2}\right)
   &=k_i^{(3)}(X)
     :=i^2X-i^2H_{i-1}^{(3)}-\frac1{2i}-\frac12.
       \label{eq:z3-functional}
\end{align}
Polynomial division and partial fractions make this definition
unambiguous. Its values are affine in $X$, and evaluation at $X=\zeta(3)$
is integration against $\mu_3$. In particular, no other odd zeta value
enters this functional.

\subsection{The Gram determinant}

For $j\ge0$ put
\[
 P_j(z)=\frac{\prod_{a=0}^{j-1}(z-a^2)}{(2j)!},\qquad
 B_j(t)=P_j(-t),\qquad D_N(t)=\prod_{i=1}^N(t+i^2),
\]
where an empty product is $1$. Define
\begin{equation}\label{eq:z3-pencil}
 A_N^{(3)}(X)
 =(N!)^2\left[
 \mathcal J_{3,X}\left(\frac{B_jB_k}{D_N}\right)
 \right]_{0\le j,k<N},\qquad
 F_N^{(3)}(X)=\det A_N^{(3)}(X).
\end{equation}
Every entry belongs to $\Q[X]$ and has degree at most one, so
$\deg F_N^{(3)}\le N$. In fact the degree is exactly $N$.
Set $V_{ji}=B_j(-i^2)$ for $0\le j<N$, $1\le i\le N$.
Partial fractions show that the coefficient matrix of $X$ is
\[
 (N!)^2 V\,\operatorname{diag}\left(
       \frac{i^2}{D_N'(-i^2)}\right)_{i=1}^N V^{\mathsf T}.
\]
Evaluation at $N$ distinct nodes is an isomorphism on the polynomials
of degree below $N$, so $V$ is invertible. Every diagonal factor is
nonzero, proving the degree assertion. At $X=\zeta(3)$ this is a Gram matrix for
the linearly independent polynomials $B_0,\ldots,B_{N-1}$ against the
positive measure $(N!)^2\dd\mu_3/D_N$. A nonzero polynomial has only
finitely many zeros, whereas this measure has positive density on
$(0,\infty)$. Thus
\begin{equation}\label{eq:z3-positive}
 F_N^{(3)}(\zeta(3))>0.
\end{equation}

For comparison with the classical linear forms, the first determinant is
$F_1^{(3)}(X)=X-1$, and the second pencil and determinant are
\begin{equation}\label{eq:z3-small-example}
 A_2^{(3)}(X)=
 \begin{pmatrix}
 5-4X&12-10X\\
 12-10X&101/4-21X
 \end{pmatrix},\qquad
 4F_2^{(3)}(X)=-64X^2+136X-71.
\end{equation}
The off-diagonal entry is $-2(5X-6)$, containing the familiar first
Ap\'ery linear form. The determinant itself is already quadratic.
This distinguishes the displayed sequence from the usual sequence of
linear forms without asserting that the underlying approximation
constructions are unrelated.

\subsection{An integer coefficient-clearing factor}

For a prime $p$, let $v_p$ on a rational polynomial mean the minimum
valuation of its coefficients. Write
$\mathcal O_p=\{a/b\in\Q:p\nmid b\}$.
We shall construct positive integers $M_N$ such that
\begin{equation}\label{eq:z3-clearing-goal}
 M_NF_N^{(3)}\in\Z[X],\qquad
 \limsup_{N\to\infty}\frac{\log M_N}{N^2}
 \le2+2\log2.
\end{equation}
Primitivity is unnecessary for this conclusion.

\begin{lemma}\label{lem:z3-local-rank}
Suppose $p>3$ and $p^2>4N+20$. Write
\[
 N=qp+\sigma,\quad 0\le\sigma<p,\quad
 \epsilon=\floor{2\sigma/p},\quad
 D_p=\floor{(2N-2)/p},\quad
 w_p=\begin{cases}
 \sigma,&\epsilon=0,\\
 \sigma-(p-1)/2,&\epsilon=1.
 \end{cases}
\]
Then
\begin{equation}\label{eq:z3-local-rank-bound}
 v_p(F_N^{(3)})\ge
 -(2+\epsilon)N-\min\{N,(D_p+1)w_p\}.
\end{equation}
\end{lemma}
\begin{proof}
The proof of \cref{prop:local-kernel} applies with the following explicit
changes. Let $V$ be the $\mathcal O_p$-span of $P_j(x^2)$, $0\le j<N$.
By \cref{lem:newton-tools}, its elements are integer-valued at $p$, and
the reduction of $f(c+pY)$ has degree at most $D_p$.
A reflected nonzero residue class contains the indices
$i\equiv\pm c\pmod p$, where $1\le c\le(p-1)/2$.
Its largest possible size is $2q+1+\epsilon$, attained by exactly $w_p$
classes, with an empty collection allowed. Restriction to these classes
defines a map into $w_p$ copies of $\Fp[Y]_{\le D_p}$.

Let $\Bfun$ be the Bernoulli evaluation map
$\Bfun(x^j)=\mathrm B_j$. In place of three derivatives in
\eqref{eq:bernoulli-moments}, the present polynomial moments satisfy
\begin{equation}\label{eq:z3-bernoulli-functional}
 \mathcal J_{3,X}(q(t))
 =\frac12\Bfun\left[\frac{\mathrm d}{\mathrm d x}
       \{x^3q(-x^2)\}\right].
\end{equation}
Expand a product $fg$ in the central-factorial basis and truncate at
index $N$, as in \eqref{eq:normalized-local-quotient}. On each residue
class, monic division shows that the normalized polynomial quotient is
integer-valued after multiplication by $p^{1+\epsilon}$. Its degree,
after multiplication by $x^3$, plus one is less than $p^2$.
One derivative and Bernoulli evaluation each lose at most one power
of $p$, by \eqref{eq:bernoulli-newton}--\eqref{eq:newton-derivative}.
Thus its contribution has valuation at least $-(3+\epsilon)$.
If one argument is in the restriction kernel, the numerator gains one
power of $p$ on every maximal class, exactly as in the proof of
\cref{prop:local-kernel}; the resulting bound is $-(2+\epsilon)$.
The constant quotient of the retained $P_N$ term costs at most
$\epsilon$, since the zeroth moment is $1/4$.

For the proper fractions, $k_i^{(3)}(X)$ has valuation at least $-3$
when $p\nmid i$. The normalized residue factor
$(N!)^2/D_N'(-i^2)$ has valuation $2q-h_c+1$, where $h_c$ is the
size of its reflected class. Only maximal classes can reach
$-(3+\epsilon)$, and a restriction-kernel argument again gains one
power. If $p\mid i$, the normalized residue factor has valuation $2$
and $k_i^{(3)}$ has valuation at least $-1$, so these contributions
are integral. Finally, the coefficient of $X$ has valuation at least
$-\epsilon$.

It follows that $p^{3+\epsilon}A_N^{(3)}$ is integral, its reduction is
constant in $X$, and its reduced rank is at most $(D_p+1)w_p$.
If that rank is $k$, row reduction supplies a factor $p$ in each of
$N-k$ rows. Removing the scale $p^{3+\epsilon}$ proves
\eqref{eq:z3-local-rank-bound}.
\end{proof}

Small primes can be cleared uniformly. Put
$L_p=\floor{\log_p(4N+4)}$. The proof of
\cref{prop:small-primes}, with cubic harmonic sums and the single
derivative in \eqref{eq:z3-bernoulli-functional}, gives the conservative
entry bound
\[
 v_p(A_N^{(3)}(X)_{jk})\ge-9L_p-v_p(96).
\]
Indeed, the proper-fraction estimate loses at most $3L_p$ in the
harmonic sum instead of $5L_p$, and the polynomial-quotient estimate
uses one derivative instead of three. All other factorial and
product-basis bounds in that proof remain valid. With $Y=N^{9/10}$,
the integer
\[
 E_N=96\prod_{p\le Y}p^{9L_p}
\]
clears these primes entrywise, and
$\log(E_N^N)=O(N^{19/10}\log N)=o(N^2)$.

For $p>N$, pass to the monomial matrix
$[\mathcal J_{3,X}(t^{j+k}/D_N)]_{j,k<N}$. It is integral at $p$.
To check this assertion, its polynomial quotients have degree at most
$N-2$. In \eqref{eq:z3-moments}, the first possible Bernoulli pole,
$2j+2=p-1$, is canceled by $2j+3=p$; the next possible pole has
$j=p-2>N-2$. All required moments are therefore integral.
Every residue node is below $p$. The reflected-pair cancellation in
\cref{lem:large-prime-integrality} still holds, now using
$\sum_{a=1}^{p-1}a^{-3}\equiv0\pmod p$, which yields
$k_{p-i}^{(3)}\equiv k_i^{(3)}\pmod p$. Singletons have unit
derivative denominators. The change of basis is exact:
\[
 F_N^{(3)}(X)=
 \frac{(N!)^{2N}}{\prod_{j=0}^{N-1}((2j)!)^2}
 \det\left[\mathcal J_{3,X}\left(\frac{t^{j+k}}{D_N}\right)\right]_{j,k<N}.
\]
Thus, for $N<p\le2N$, a clearing exponent is
$2\sum_{j=0}^{N-1}\floor{2j/p}$, and no prime beyond $2N$ is needed.
For sufficiently large $N$, we have $Y>3$ and $Y^2>4N+20$.
With $\epsilon_p,D_p,w_p$ as in \cref{lem:z3-local-rank}, set
\[
 M_N=E_N^N
 \prod_{Y<p\le N}
 p^{(2+\epsilon_p)N+\min\{N,(D_p+1)w_p\}}
 \prod_{N<p\le2N}p^{\,2\sum_{j=0}^{N-1}\floor{2j/p}}.
\]
All exponents are nonnegative integers, and the preceding estimates
prove $M_NF_N^{(3)}\in\Z[X]$. For the finitely many remaining $N$,
take any positive integer clearing the coefficients of $F_N^{(3)}$.

We evaluate its rate. Away from the floor boundaries put $\beta=p/N$,
$q=\floor{1/\beta}$, and
$\epsilon=\floor{2\{1/\beta\}}$. On $0<\beta<1$, the limiting
clearing profile is
\[
 e_3(\beta)=2+\epsilon+
 \min\left\{1,
 (2q+\epsilon+1)\left(1-(q+\epsilon/2)\beta\right)\right\}.
\]
The expression inside the minimum ranges from $0$ to $1$ on each
half-period, so explicitly
\begin{equation}\label{eq:z3-arithmetic-profile}
 e_3(\beta)=
 \begin{cases}
 2q+5-(q+1)(2q+1)\beta,
   &\dfrac1{q+1}<\beta<\dfrac2{2q+1},\\[2mm]
 2q+3-q(2q+1)\beta,
   &\dfrac2{2q+1}<\beta<\dfrac1q.
 \end{cases}
\end{equation}
Its integral over the $q$th full period is
\[
 \int_{1/(q+1)}^{1/q}e_3(\beta)\,\dd\beta
 =\frac{12q+5}{2q(q+1)(2q+1)}
 =\frac5{2q}-\frac7{2(q+1)}+\frac2{2q+1}.
\]
Summing to $Q$ gives
$2(H_{2Q+1}-H_Q)+3/2-7/(2Q+2)$, where
$H_j=\sum_{i=1}^j i^{-1}$. Consequently
\[
 \int_0^1e_3(\beta)\,\dd\beta=\frac32+2\log2.
\]
Above $1$, the factorial charge has profile $2-\beta$ on $(1,2)$,
whose integral is $1/2$. The prime-number-theorem passage in
\cref{sec:arithmetic} therefore proves \eqref{eq:z3-clearing-goal}.
For precision, $e_3\le4$ controls the interval near zero; after fixing
a positive lower cutoff there are finitely many floor boundaries, each
with only $O(1)$ exceptional integer primes and total cost
$O(N\log N)$. Small primes have already cost $o(N^2)$. Thus no
prime-number theorem for shrinking intervals is required.

\subsection{The analytic saving at the unshifted window}

Let $\widehat h_{k,N}$ be the least squared norm of a real monic
polynomial of degree $k$ for $\dd\mu_3/D_N$. Positive definiteness
gives the minimizer, and the triangular change to monic orthogonal
polynomials gives
\begin{equation}\label{eq:z3-norm-product}
 F_N^{(3)}(\zeta(3))
 =\prod_{k=0}^{N-1}\frac{(N!)^2}{((2k)!)^2}\widehat h_{k,N}.
\end{equation}
This is the $m=0$ version of \cref{lem:norm-product}.

We specify the potential used to bound these norms. Put
\[
 I_1(y)=\int_0^1\log(y^2+t^2)\,\dd t,\qquad
 V_0(y)=2\pi y+I_1(y).
\]
For $b,t>0$, let $\lambda_b$ and $\nu_{b,t}$ be the symmetric measures
on $[-b,b]$ with respective densities
\[
 \frac1\pi\arcosh\frac b{|x|},\qquad
 \frac{t\sqrt{b^2-x^2}}
 {2\pi\sqrt{t^2+b^2}(x^2+t^2)}.
\]
Here $\arcosh z=\log(z+\sqrt{z^2-1})$ for $z\ge1$.
The positive measure
$\eta_b=\lambda_b+\int_0^1\nu_{b,t}\,\dd t$ has mass
\begin{equation}\label{eq:z3-mass}
 u(b)=\frac{3b+1-\sqrt{1+b^2}}2,
 \qquad
 b(u)=\frac{6u-3+\sqrt{4u^2-4u+9}}8.
\end{equation}
In particular $b(1)=3/4$. Write
$U_b(y)=\int_{-b}^b\log|y-x|\,\dd\eta_b(x)$.
The Cauchy-transform computation \eqref{eq:cauchy-nu} and its
integration in \cref{prop:potential-inequality} give
\begin{align}
 4U_b(y)-V_0(y)&=\ell(u) &&(0\le y\le b),\notag\\
 4U_b(y)-V_0(y)&\le\ell(u) &&(y\ge0),\notag\\
 \ell(u)&=2+4u\left(\log\frac b2-1\right)
              -2\log\frac{1+\sqrt{1+b^2}}2,\quad b=b(u).
 \label{eq:z3-potential}
\end{align}
These formulas follow directly with the term $\int_0^s\nu_{b,t}\,\dd t$
absent. In particular, $\eta_b\le(3/2)\lambda_b$,
$u(b)\sim3b/2$ at zero, and $\ell(0)=2=-V_0(0)$.
They do not use the different small-$b$ expansion valid at fixed $s>0$.

\begin{proposition}\label{prop:z3-analytic}
Define $g(0)=0$ and, for $0<u\le1$,
\begin{equation}\label{eq:z3-g}
 g(u)=4u\log\frac{4u}{b(u)}
       +2\log\frac{1+\sqrt{1+b(u)^2}}2.
\end{equation}
Then
\begin{equation}\label{eq:z3-analytic-bound}
 \limsup_{N\to\infty}\frac{\log F_N^{(3)}(\zeta(3))}{N^2}
 \le-\int_0^1g(u)\,\dd u<-5\log2.
\end{equation}
\end{proposition}
\begin{proof}
Fix $\varepsilon>0$. The quantile construction of
\cref{lem:quantile-bound} applies uniformly for $\varepsilon\le u\le1$:
$b(u)$ stays in a compact positive interval and
$\eta_b\le(3/2)\lambda_b$ gives the same folded-density bound
$Cv^{-1/2}(1+|\log v|)$. It supplies monic trial polynomials of degree
$k$ whose squared values at $N^2y^2$ are at most
$N^{4k+C_\varepsilon}e^{4NU_b(y)}$, with $u=k/N$.
Also $D_N(N^2y^2)\ge N^{2N}e^{NI_1(y)}$.
Together with \eqref{eq:z3-density-bound} and
\eqref{eq:z3-potential}, this gives
\[
 \log\widehat h_{k,N}
 \le(4k-2N)\log N+N\ell(k/N)+O_\varepsilon(\log N)
 \quad(\varepsilon N\le k<N).
\]
The tail of $4U_b-V_0$ decreases at a uniform linear rate outside a
fixed compact interval, as in \eqref{eq:tail-derivative}, so the
pointwise estimate may be integrated. Replacing the old density's
polynomial power $5$ by $3$ only changes the logarithmic error.
Stirling's formula in \eqref{eq:z3-norm-product} then bounds the
logarithm of its $k$th factor by $-Ng(k/N)+O_\varepsilon(\log N)$.

For $0\le k<\varepsilon N$, use $B_k$ itself as a trial polynomial.
The corresponding factor is at most
\[
 (N!)^2\int\frac{B_k(t)^2}{D_N(t)}\,\dd\mu_3(t)
 \le\int B_k(t)^2\,\dd\mu_3(t)
 \le Ce^{Ck}(1+k)^3.
\]
The last inequality follows from
$|B_k(x^2)|\le(x+k)^{2k}/(2k)!$ and the gamma integral in
\eqref{eq:z3-density-bound}; at $k=0$ the integral is $1/4$.
These initial rows contribute $O(\varepsilon^2N^2)+O(N\log N)$.
Letting $N\to\infty$ and then $\varepsilon\downarrow0$ proves the
first inequality in \eqref{eq:z3-analytic-bound}. Continuity at zero
follows from $g(u)=4u\log6+O(u^2)$.

Finally, differentiation gives $g'(u)=4\log(4u/b)$. Since
\[
 u-bu'(b)=\frac{1-1/\sqrt{1+b^2}}2>0\qquad(b>0),
\]
$g$ is strictly concave. Its endpoint values are
$g(0)=0$ and $g(1)=10\log2$, so its graph lies strictly above the
endpoint chord on $(0,1)$. Integrating proves
$\int_0^1g(u)\,\dd u>5\log2$.
\end{proof}

The exact value is also elementary:
\begin{equation}\label{eq:z3-exact-analytic}
 \int_0^1g(u)\,\dd u=\frac{13}{2}\log2-1.
\end{equation}
For verification, set $z=\sqrt{1+b^2}+b$, so $1\le z\le2$ and
$u=(z+1-2/z)/2$. Put
$L(z)=\log(4(z+2)/(z+1))$ and
$M(z)=\log((z+1)^2/(4z))$. Then $g=4uL+2M$ and $dg/du=4L$.
Integration by parts gives
\[
 \int_0^1g\,\dd u
 =2L(2)+2M(2)
 -2\int_1^2\frac{u(z)^2}{(z+1)(z+2)}\,\dd z.
\]
Here $L(2)+M(2)=\log6$, and
\[
 \frac{u(z)^2}{(z+1)(z+2)}
 =\frac14\left(1-\frac5z+\frac2{z^2}+\frac4{z+1}\right).
\]
The last integral is $1/2+\log3-(9/4)\log2$, proving
\eqref{eq:z3-exact-analytic}. The endpoint term at $u=0$ vanishes
because $L(z)\to\log6$ there.

\subsection{The contradiction and the role of the warm-up}

\begin{proof}[Proof of \cref{thm:z3-warmup}]
Suppose $\zeta(3)=a/b$ with $a,b\in\Z$ and $b>0$.
By \eqref{eq:z3-clearing-goal}, the degree bound, and
\eqref{eq:z3-positive},
\[
 Z_N=b^NM_NF_N^{(3)}(a/b)\in\Z_{>0}.
\]
But \eqref{eq:z3-clearing-goal} and \eqref{eq:z3-analytic-bound} give
\[
 \limsup_{N\to\infty}\frac{\log Z_N}{N^2}
 \le2-3\log2<0.
\]
Indeed,
$\log2=2\int_0^{1/3}(1-t^2)^{-1}\,\dd t>2/3$.
Thus $Z_N\to0$, contradicting $Z_N\ge1$.
\end{proof}

Using the exact integral instead of the chord bound gives the stronger
explicit margin
\[
 \left(\frac{13}{2}\log2-1\right)-(2+2\log2)
 =\frac92\log2-3>0.
\]
This is a comparison of proven sufficient estimates, not a claim that
the coefficient denominator or the determinant has a matching
asymptotic lower bound. The subsequent $\zeta(5)$ argument keeps the
same positive-functional and Gram-determinant structure. Its larger
harmonic-denominator cost is overcome here by a shifted window,
additional local divisibility, and a more delicate final comparison. The present
case therefore exhibits the mechanism before those refinements are
introduced.

% END JOURNAL ADDITION zeta3

\medskip\noindent\textbf{The $\zeta(5)$ proof.} The parameter conventions in the following sections apply through Section~\ref{sec:contradiction}; Section~\ref{sec:higher-odd} subsequently reintroduces the general odd target.

% BEGIN SHARED CORE 02_plan
\section{Plan of the proof and notation}\label{sec:plan}

We give a candidate proof of the following statement, independently of the
arguments cited in the historical discussion for the same conclusion.
Here independence specifies the logical dependencies: no theorem or lemma
from Fauzan's proof or its subsequent adaptations is used below.

\begin{theorem}\label{thm:main}
The real number
\[
  \zeta(5)=\sum_{n=1}^{\infty}\frac1{n^5}
\]
is irrational.
\end{theorem}

The proof has two parts. A positive measure produces a determinant
$F_N(X)\in\Q[X]$ that is strictly positive at $X=\zeta(5)$. Its entries are
affine functions of $X$, and $\deg F_N\le N$. We estimate both the denominators
of its coefficients and its value at $\zeta(5)$ on the scale $N^2$. The
arithmetic estimate costs less than $5.4101N^2$ in the logarithm, while the
analytic estimate saves more than $5.4133N^2$. If $\zeta(5)$ were rational,
clearing its denominator would cost only $O(N)$. The resulting positive
integer would tend to zero.

\subsection{Parameters and conventions}
Throughout, $N$ tends to infinity through positive multiples of $20$, and
\begin{equation}\label{eq:fixed-parameters}
 s=\frac1{20},\qquad \rho=\frac54,\qquad
 \delta=\rho-s=\frac65,\qquad
 m=sN,\qquad R=\rho N,\qquad r=R-m=\delta N.
\end{equation}
The letter $r$ denotes an integer matrix size; $\delta$ denotes its ratio to
$N$. For an integer $j\ge0$, empty products being $1$, put
\begin{equation}\label{eq:basic-polynomials}
 P_j(z)=\frac{\prod_{a=0}^{j-1}(z-a^2)}{(2j)!},\qquad
 B_j(t)=P_j(-t),\qquad
 D_N(t)=\prod_{i=1}^{N}(t+i^2).
\end{equation}
The \emph{principal factorial window} is simply the consecutive collection
$B_m,\ldots,B_{R-1}$. This phrase describes an index range and its factorial
normalization; it introduces no additional mathematical structure.

All logarithms are natural. We write $x_+=\max(x,0)$ and use the usual floor
and ceiling functions. The notation $O(\cdot)$ and $o(\cdot)$ refers to the
limit just specified; a claimed uniformity in other variables is stated
explicitly. Constants in intermediate estimates may depend on the fixed
parameters in \eqref{eq:fixed-parameters}. An integral against a measure means
a Lebesgue integral, or its improper version when its convergence has been
established.

For a prime $p$, $v_p(a/b)$ is the exponent of $p$ in a nonzero rational
number $a/b$, and $v_p(0)=+\infty$. For a rational polynomial, vector, or
matrix, $v_p$ means the minimum of the valuations of its coefficients or
entries. Thus a polynomial inequality $v_p(G)\ge -e$ is
\emph{coefficientwise}: it concerns every coefficient, not merely the value
of $G$ at one point. We use the localization
\[
 \mathcal O_p=\Z_{(p)}
 =\{a/b:a,b\in\Z,\ b\ne0,\ p\nmid b\},
\]
whose reduction modulo $p$ is the field $\Fp=\Z/p\Z$. A \emph{unit} of
$\mathcal O_p$ is an element with valuation zero. These definitions suffice;
no completion of $\Q$ is required.

The Bernoulli numbers $\mathrm B_n$ are defined by the convergent power
series near zero, equivalently by its identity of formal series,
\[
 \frac{z}{e^z-1}=\sum_{n=0}^{\infty}\mathrm B_n\frac{z^n}{n!}.
\]
The upright symbol $\mathrm B_n$ distinguishes these numbers from the
polynomials $B_j(t)$. We write
$H_n^{(5)}=\sum_{i=1}^{n}i^{-5}$, with $H_0^{(5)}=0$.
The linear functional $\J_X$ and the positive measure underlying it are
defined in \cref{sec:functional}.

\subsection{The two estimates to be proved}
For a nonzero polynomial $G(X)=\sum_jg_jX^j\in\Q[X]$, its
\emph{primitive positive normalizer} is the unique positive rational number
$Q(G)$ such that $Q(G)G$ has integer coefficients with greatest common divisor
$1$. Explicitly,
\[
 Q(G)=\prod_p p^{-\min_j v_p(g_j)}.
\]
Only finitely many factors differ from $1$. The number $Q(G)$ can be less
than $1$ when all coefficients of $G$ share integer factors. Keeping these
factors is essential to the estimate. Write $Q_N=Q(F_N)$.

The two quantitative conclusions are
\begin{align}
 \limsup_{N\to\infty}\frac{\log Q_N}{N^2}
    &\le C_{\arith}<\frac{54101}{10000},\label{eq:plan-arith}\\
 \limsup_{N\to\infty}\frac{\log F_N(\zeta(5))}{N^2}
    &\le -C_{\an},\qquad C_{\an}>\frac{54133}{10000}.
       \label{eq:plan-analytic}
\end{align}
The exact rational certificate in \cref{sec:certificate} proves the displayed
strict comparisons without floating-point arithmetic. It does not replace
the proofs of the estimates that lead to those comparisons.

Here is the organization. Sections~\ref{sec:functional}--\ref{sec:determinant}
construct the positive determinant.
Sections~\ref{sec:integer-valued}--\ref{sec:root-gain} establish local
arithmetic estimates, including the improvement from the common initial
roots of the window.
Sections~\ref{sec:potential}--\ref{sec:discretization} give an explicit
logarithmic-potential construction and convert it into polynomial norm
bounds.
Section~\ref{sec:arithmetic} sums the local estimates over primes.
Sections~\ref{sec:certificate}--\ref{sec:contradiction} certify the numerical
gap and finish the argument.

For clarity, the terms \emph{local restriction kernel}, \emph{high atom},
\emph{quotient support}, and \emph{arithmetic profile} are descriptive names
defined at their first use below. No familiarity with a separate framework
bearing these names is assumed. The historical formalizations in
\cref{sec:history} concern their own stated arguments; they are not formal
certifications of the proof developed here.

% END SHARED CORE 02_plan

% BEGIN SHARED CORE 03_functional
\section{A positive measure with rational moments and affine Stieltjes values}\label{sec:functional}

We construct one positive measure that will serve both the arithmetic and
the analytic parts of the argument. Its polynomial moments are rational,
whereas its values on the simple fractions $(t+i^2)^{-1}$ contain exactly
one common undetermined number, namely $\zeta(5)$.

\begin{definition}\label{def:physical-measure}
For $x>0$, put $f(x)=(e^{2\pi x}-1)^{-1}$ and
\[
 w(x)=\frac{x^5}{12}f^{(4)}(x).
\]
The base measure $\mu$ on $[0,\infty)$ is the pushforward of
$w(x)\,\dd x$ under $x\mapsto x^2$. Explicitly, for every nonnegative
measurable function $g$, this means
\begin{equation}\label{eq:physical-measure}
 \int_{[0,\infty)}g(t)\,\dd\mu(t)
 =\frac1{12}\int_0^\infty g(x^2)x^5f^{(4)}(x)\,\dd x.
\end{equation}
Thus an integral written in the variable $t$ uses $\mu$, and the same
integral written in the variable $x$ uses the displayed density $w$.
\end{definition}

\begin{lemma}\label{lem:positive-measure}
The density $w$ is strictly positive on $(0,\infty)$, extends continuously
to $w(0)=1/\pi$, and satisfies
$w(x)\le C(1+x)^5e^{-2\pi x}$ for $x\ge0$, with an absolute constant $C$.
Every moment is finite and rational. More precisely, for every integer
$k\ge0$,
\begin{equation}\label{eq:rational-moments}
 \mu_k:=\int t^k\,\dd\mu(t)
 =\frac{(2k+5)!}{12(2\pi)^{2k+2}}\zeta(2k+2)
 =\frac{(-1)^k}{24}(2k+3)(2k+4)(2k+5)\mathrm B_{2k+2}.
\end{equation}
\end{lemma}
\begin{proof}
On each interval $x\ge a>0$, termwise differentiation of the locally
uniformly convergent exponential series gives
\[
 f^{(4)}(x)=(2\pi)^4\sum_{n=1}^{\infty}n^4e^{-2\pi nx}>0.
\]
The Laurent expansion $f(x)=(2\pi x)^{-1}-1/2+O(x)$ at zero, with its
convergent power series remainder, gives $w(0)=1/\pi$. For $x\ge1$, the
exponential series bounds $f^{(4)}(x)$ by a constant times $e^{-2\pi x}$;
continuity supplies the asserted bound on the remaining compact interval.
Tonelli's theorem and the elementary gamma integral now give
\[
 \mu_k=\frac{(2\pi)^4}{12}
 \sum_{n\ge1}n^4\frac{(2k+5)!}{(2\pi n)^{2k+6}},
\]
which is the first identity. The second is Euler's even-zeta identity
\cite[\S25.6.2]{dlmf}, with the Bernoulli convention of
\cref{sec:plan}. It also follows by comparing the Taylor expansion with
the partial fractions of $\coth\pi x$. No statement concerning odd zeta
values is used in this step.
\end{proof}

\begin{lemma}\label{lem:stieltjes}
For $y>0$, define the Hurwitz zeta value by the convergent series
$\zeta(5,y)=\sum_{n=0}^{\infty}(n+y)^{-5}$. Then
\begin{equation}\label{eq:stieltjes-hurwitz}
 \int\frac{\dd\mu(t)}{t+y^2}
 =y^4\zeta(5,y)-\frac1{2y}-\frac14.
\end{equation}
In particular, for every positive integer $i$,
\begin{equation}\label{eq:stieltjes-integer}
 \int\frac{\dd\mu(t)}{t+i^2}
 =i^4\zeta(5)-i^4H_{i-1}^{(5)}-\frac1{2i}-\frac14.
\end{equation}
\end{lemma}
\begin{proof}
Put $K(y)=\int_0^\infty xf(x)/(x^2+y^2)\,\dd x$. The elementary Binet
integral for the digamma function is
\begin{equation}\label{eq:binet-digamma}
 2K(y)=\log y-\frac1{2y}-\psi(y),
 \qquad
 \psi(y)=-\gamma+\sum_{n=0}^{\infty}
 \left(\frac1{n+1}-\frac1{n+y}\right).
\end{equation}
Here $\gamma=\lim_{M\to\infty}(\sum_{n=1}^M n^{-1}-\log M)$.
For completeness, the partial fraction identity
$\coth\pi x=(\pi x)^{-1}+(2x/\pi)\sum_{n\ge1}(x^2+n^2)^{-1}$
and the Laplace sine transform give
\[
 2K(y)=\int_0^\infty e^{-yt}
 \left(\frac1{e^t-1}-\frac1t+\frac12\right)\dd t.
\]
Differentiation of this last integral gives
$-\sum_{n\ge1}(n+y)^{-2}+y^{-1}-(2y^2)^{-1}$.
This equals the derivative of the right side of
\eqref{eq:binet-digamma}, by differentiating the defining series for
$\psi$. Both sides tend to zero as $y\to\infty$; for $\psi$ this follows
by comparing its defining sum with the integral of $1/(y+t)$.
This proves the identity, also recorded in \cite[\S5.9.15]{dlmf}.
The integrals and their derivatives converge uniformly when $y$ stays in
a compact subinterval of $(0,\infty)$.

Let $w_y(x)=x^5/(x^2+y^2)$. In four integrations by parts, all the boundary
products vanish: at zero each product
$w_y^{(a)}(x)f^{(3-a)}(x)$, $0\le a\le3$, is $O(x)$, and at infinity
the derivatives of $f$ decay exponentially. Therefore
\[
 \int\frac{\dd\mu(t)}{t+y^2}
 =\frac1{12}\int_0^\infty w_y^{(4)}(x)f(x)\,\dd x.
\]
The decomposition
$w_y=x^3-y^2x+y^4x/(x^2+y^2)$ and harmonicity of
$x/(x^2+y^2)$ imply
$w_y^{(4)}=y^4\partial_y^4[x/(x^2+y^2)]$.
Consequently the last integral is $y^4K^{(4)}(y)/12$.
Differentiating \eqref{eq:binet-digamma} four times gives
\[
 K^{(4)}(y)=-\frac3{y^4}-\frac6{y^5}+12\zeta(5,y),
\]
which proves \eqref{eq:stieltjes-hurwitz}. Splitting the defining zeta
series at $i$ proves \eqref{eq:stieltjes-integer}.
\end{proof}

\begin{definition}\label{def:formal-functional}
Let $\mathcal R$ be the $\Q$-vector space of rational functions with poles
only at finitely many points $-i^2$, $i\ge1$, all of order at most one.
Every element has a unique expression
$q(t)+\sum_i a_i/(t+i^2)$, with $q\in\Q[t]$ and $a_i\in\Q$.
Define the $\Q$-linear functional $\J_X:\mathcal R\to\Q[X]$ by
\begin{equation}\label{eq:formal-functional}
 \J_X(t^k)=\mu_k,\qquad
 \J_X\left(\frac1{t+i^2}\right)=k_i(X)
 :=i^4X-i^4H_{i-1}^{(5)}-\frac1{2i}-\frac14.
\end{equation}
\end{definition}

Polynomial division and the uniqueness of partial fractions make this
definition unambiguous. Its values have degree at most one in $X$.
The functional is linear, not multiplicative. By the preceding lemmas,
$\J_{\zeta(5)}(g)=\int g\,\dd\mu$ for every $g\in\mathcal R$; the integral
converges because all poles lie on the negative real axis and $\mu$ has
all moments.

% END SHARED CORE 03_functional

% BEGIN SHARED CORE 04_determinant
\section{The determinant and its primitive coefficient normalizer}\label{sec:determinant}

Retain the parameters and polynomials in
\eqref{eq:fixed-parameters}--\eqref{eq:basic-polynomials}. It is useful
to distinguish the full matrix from its principal window. Define
\begin{equation}\label{eq:full-pencil}
 A_{N,R}(X)=(N!)^2
 \left[\J_X\left(\frac{B_j(t)B_k(t)}{D_N(t)}\right)\right]_{0\le j,k<R},
\end{equation}
and let
\begin{equation}\label{eq:window-pencil}
 A_N(X)=\bigl[A_{N,R}(X)_{jk}\bigr]_{m\le j,k<R},
 \qquad F_N(X)=\det A_N(X).
\end{equation}
An affine matrix in the variable $X$, such as this one, is sometimes
called a matrix pencil. Every rational function in these definitions
belongs to $\mathcal R$, including when the numerator degree exceeds
the denominator degree. Its polynomial quotient is part of the
definition of $\J_X$ and will be retained in every arithmetic estimate.

\begin{proposition}\label{prop:gram-pencil}
The polynomial $F_N$ is nonzero, belongs to $\Q[X]$, and has degree at
most $N-m+1$, hence at most $N$. Moreover,
$F_N(\zeta(5))>0$.
\end{proposition}
\begin{proof}
At $X=\zeta(5)$ the matrix is the Gram matrix, multiplied by $(N!)^2$,
of $B_m,\ldots,B_{R-1}$ for the positive measure $\dd\mu/D_N$.
These polynomials have distinct degrees and are linearly independent.
For a nonzero real vector $(a_j)$, their linear combination is a
nonzero polynomial, so its square has strictly positive integral against
this measure: its zero set is finite and the measure has a positive
density on $(0,\infty)$. Thus the matrix is positive definite and its
determinant is positive.

The coefficient of $X$ in an entry is obtained from the residues alone.
Writing $t_i=-i^2$, it is
\[
 (N!)^2\sum_{i=1}^N
 \frac{i^4}{D_N'(t_i)}B_j(t_i)B_k(t_i).
\]
Each summand is a matrix of rank at most one. If $i<m$, then
$B_j(t_i)=P_j(i^2)=0$ for every $j\ge m$. Hence the coefficient matrix
has rank at most $N-m+1$. In the column expansion of
$\det(A_0+XA_1)$, a term with more than $\rank A_1$ columns taken
from $A_1$ is zero. This proves the degree bound. Rationality follows
from \cref{def:formal-functional}.
\end{proof}

\begin{lemma}\label{lem:primitive-normalizer}
Write $F_N(X)=\sum_j c_{N,j}X^j$ and put
\begin{equation}\label{eq:primitive-normalizer}
 \nu_{N,p}=\min_{c_{N,j}\ne0}v_p(c_{N,j}),\qquad
 Q_N=\prod_p p^{-\nu_{N,p}}.
\end{equation}
Then $Q_NF_N\in\Z[X]$ has relatively prime coefficients. If
$\widetilde Q_N>0$ is rational and
$\widetilde Q_NF_N\in\Z[X]$, then
\begin{equation}\label{eq:minimal-normalizer}
 \widetilde Q_N/Q_N\in\Z_{>0},\qquad Q_N\le\widetilde Q_N.
\end{equation}
\end{lemma}
\begin{proof}
Only finitely many primes divide a numerator or denominator of the
nonzero coefficients. At each prime, the minimum coefficient valuation
of $Q_NF_N$ is zero, proving both integrality and primitivity.
The integrality of $\widetilde Q_NF_N$ says
$v_p(\widetilde Q_N)+\nu_{N,p}\ge0$ for every prime. These are exactly
the assertions that the rational number $\widetilde Q_N/Q_N$ has
nonnegative valuation at every prime, and hence is an integer.
\end{proof}

The normalizer can contain negative prime powers, because common
coefficient factors are retained. Accordingly, a local clearing exponent
below is allowed to be negative; it then asserts common divisibility
rather than a denominator.

\begin{corollary}\label{cor:closing-criterion}
If $\zeta(5)=a/b$ for integers $a,b$ with $b>0$, then
\begin{equation}\label{eq:positive-integer}
 b^NQ_NF_N(a/b)\in\Z_{>0}.
\end{equation}
Thus any estimates
$\log Q_N\le (c+o(1))N^2$ and
$\log F_N(\zeta(5))\le-(c+\eta+o(1))N^2$, with $\eta>0$,
give a contradiction.
\end{corollary}
\begin{proof}
The degree bound makes $b^N(Q_NF_N)(a/b)$ an integer, and positivity
makes it a positive one. Under the two estimates its logarithm is at
most $N\log b-\eta N^2+o(N^2)$, which tends to $-\infty$.
\end{proof}

\subsection{A second basis for the arithmetic estimates}
Define the common root polynomial and an $r\times r$ matrix by
\[
 C_m(t)=\prod_{a=0}^{m-1}(-t-a^2)=(2m)!B_m(t),\qquad
 H_{kl}(X)=\J_X\left(\frac{C_m(t)^2t^{k+l}}{D_N(t)}\right),
 \quad 0\le k,l<r.
\]
The phrase \emph{common roots} refers to the fact that $C_m$ divides
every polynomial in the chosen window.

\begin{lemma}\label{lem:basis-change}
There is an exact identity in $\Q[X]$:
\begin{equation}\label{eq:monomial-change}
 F_N(X)=\frac{(N!)^{2r}}{\prod_{j=m}^{R-1}((2j)!)^2}\det H(X).
\end{equation}
\end{lemma}
\begin{proof}
The polynomial $B_{m+k}/C_m$ has degree $k$ and leading coefficient
$(-1)^k/(2m+2k)!$. The change from the polynomials
$C_m,C_mt,\ldots,C_mt^{r-1}$ to $B_m,\ldots,B_{R-1}$ is therefore
triangular, with these diagonal entries. Applying the same change to
the two arguments of the bilinear form multiplies its Gram determinant
by the square of the change-of-basis determinant. The additional
factor $(N!)^2$ in every row yields \eqref{eq:monomial-change}.
\end{proof}

We will use the factorial basis for the restriction argument and the
monomial basis for the residue argument. Equation
\eqref{eq:monomial-change} accounts exactly for the cost of switching
between them. The separate factorization into monic orthogonal norms
is given in \cref{lem:norm-product}.

% END SHARED CORE 04_determinant

% BEGIN SHARED CORE 05_integer_valued
\section{Integer-valued polynomials and the small-prime estimate}\label{sec:integer-valued}

We first record elementary local algebra that will also be used for the
large primes. Throughout this section $p$ is any prime and
$\mathcal O_p=\Z_{(p)}$, as defined in \cref{sec:plan}. Reduction sends
$a/b\in\mathcal O_p$ to $\overline a\,\overline b^{-1}\in\Fp$.
A rational polynomial is \emph{integer-valued at $p$} if its value at
every integer belongs to $\mathcal O_p$. This is weaker than belonging
to $\mathcal O_p[x]$; for example, $\binom{x}{p}$ is integer-valued at
$p$ although its leading coefficient is not in $\mathcal O_p$.

For a rational polynomial $G$ in any finite number of indeterminates,
its Gauss valuation is the least valuation of its nonzero coefficients.
It satisfies
\begin{equation}\label{eq:gauss-laws}
 v_p(GH)=v_p(G)+v_p(H),\qquad
 v_p(G+H)\ge\min\{v_p(G),v_p(H)\}.
\end{equation}
For the first identity, multiply each factor by a power of $p$ to give
Gauss valuation zero. Its reduction is then a nonzero polynomial over
$\Fp$, and the product of two such polynomials is nonzero. This argument
also proves that, if an actual polynomial quotient $G/U$ exists and
$v_p(U)=0$, division by $U$ preserves the Gauss valuation.

\begin{lemma}\label{lem:newton-tools}
The following assertions hold.
\begin{enumerate}[label=(\roman*)]
\item A polynomial $g\in\Q[x]$ of degree at most $L$ is integer-valued
at $p$ if and only if it has an expansion
$g(x)=\sum_{a=0}^L d_a\binom{x}{a}$ with $d_a\in\mathcal O_p$.
\item Put $C_0^*(x)=1$ and $C_l^*(x)=2P_l(x^2)$ for $l\ge1$.
These polynomials form an $\mathcal O_p$-basis of the even
integer-valued polynomials. In particular, when $p$ is odd,
the coefficients of every product $P_j(x^2)P_k(x^2)$ in the
$P_l(x^2)$ basis belong to $\mathcal O_p$. For $p=2$, their valuations
are at least $-2$.
\item If $g$ is integer-valued at $p$ and $\deg g=L<p^2$, then for
every integer $c$,
\begin{equation}\label{eq:local-newton}
 g(c+pY)\in\mathcal O_p[Y],\qquad
 \deg\bigl(\overline{g(c+pY)}\bigr)\le\floor{L/p}.
\end{equation}
The degree of the zero polynomial is understood as $-\infty$.
\end{enumerate}
\end{lemma}
\begin{proof}
For (i), Newton interpolation gives
$d_a=\Delta^ag(0)$, where $\Delta g(x)=g(x+1)-g(x)$.
Thus integer-valuedness implies $d_a\in\mathcal O_p$.
Conversely, $\binom{x}{a}=x(x-1)\cdots(x-a+1)/a!$ is an integer at
every integer $x$, including negative $x$, proving the reverse implication.

For (ii), direct multiplication gives, for $l\ge1$,
\begin{equation}\label{eq:central-binomial}
 2P_l(x^2)=\binom{x+l}{2l}+\binom{x+l-1}{2l}.
\end{equation}
Hence $C_l^*$ is integer-valued, vanishes at $0,\ldots,l-1$, and equals
$1$ at $l$. Successively subtracting multiples of $C_0^*,C_1^*,\ldots$
therefore expands any even integer-valued polynomial with coefficients
in $\mathcal O_p$. Once a polynomial of degree at most $2L$ vanishes at
$0,1,\ldots,L$, evenness makes it identically zero, proving existence
and uniqueness of the expansion. Since $P_l=C_l^*/2$ for $l>0$,
products of two $P$'s lose at most two powers of $2$ and no odd-prime
factor.

For (iii), it suffices to treat $\binom{x}{a}$ with $a<p^2$.
The denominator $a!$ has valuation $\floor{a/p}$. Among the numerator
factors $c+pY-j$, $0\le j<a$, at least that many have constant term
divisible by $p$; extract a factor $p$ from each of them. The remaining
coefficients are in $\mathcal O_p$. If more than $\floor{a/p}$ factors
were extracted, the reduction is zero. Otherwise only those
$\floor{a/p}$ factors can contribute $Y$ to the reduction, all other
factors reducing to nonzero constants. This proves the degree bound.
\end{proof}

\begin{definition}\label{def:bernoulli-functional}
The Bernoulli evaluation functional is the $\Q$-linear map
$\Bfun:\Q[x]\to\Q$ determined by $\Bfun(x^a)=\mathrm B_a$.
It is an algebraic functional and is not asserted to be integration
against a positive measure. A subscript, as in $\Bfun_Y$, specifies its
polynomial variable.
\end{definition}

Applying this functional to $(1+z)^x$, as a formal power series in $z$,
and using the Bernoulli generating function gives
\begin{equation}\label{eq:bernoulli-newton}
 \Bfun\binom{x}{a}=\frac{(-1)^a}{a+1}.
\end{equation}
Differentiating the same binomial generating function gives
\begin{equation}\label{eq:newton-derivative}
 \frac{\dd}{\dd x}\binom{x}{a}
 =\sum_{b=1}^{a}\frac{(-1)^{b-1}}b\binom{x}{a-b}.
\end{equation}
Consequently, on integer-valued polynomials of degree at most $L$,
one derivative loses at most $\floor{\log_p L}$ powers of $p$, and
Bernoulli evaluation loses at most $\floor{\log_p(L+1)}$ powers.
This assertion means that multiplying by the stated power of $p$
makes the result integer-valued, or in $\mathcal O_p$, respectively.
For constants the derivative is zero and no loss is needed.

We will also use two exact identities:
\begin{align}
 \J_X(q(t))
 &=\frac1{24}\Bfun\left[
       \frac{\dd^3}{\dd x^3}\{x^5q(-x^2)\}\right],
       &&q\in\Q[t],\label{eq:bernoulli-moments}\\
 \Bfun(F(x))
 &=\frac1p\sum_{c=0}^{p-1}\Bfun_Y(F(c+pY)),
       &&F\in\Q[x].\label{eq:bernoulli-distribution}
\end{align}
The first follows on each monomial from \eqref{eq:rational-moments}.
For the second, applying both sides to $e^{zx}$ as a formal generating
function gives $z/(e^z-1)$ on each side: the finite geometric sum
over $c$ cancels the factor $e^{pz}-1$.

\begin{proposition}\label{prop:small-primes}
For any prime $p$, put $L_p=\floor{\log_p(4R+4)}$. Every entry of the
full matrix in \eqref{eq:full-pencil} satisfies
\begin{equation}\label{eq:small-prime-entry}
 v_p(A_{N,R}(X)_{jk})\ge-9L_p-v_p(96).
\end{equation}
For all sufficiently large $N$, the positive integer
\begin{equation}\label{eq:small-prime-factor}
 E_N=96\prod_{p\le N^{9/10}}p^{9L_p}
\end{equation}
clears these primes in every entry, and $E_N^r$ clears them in $F_N$.
Furthermore, $\log(E_N^r)=o(N^2)$.
\end{proposition}
\begin{proof}
The product expansion in \cref{lem:newton-tools} reduces the estimate to
\[
 T_l(X)=(N!)^2\J_X\left(\frac{P_l(-t)}{D_N(t)}\right),
 \qquad 0\le l\le2R-2,
\]
with an additional loss of at most $2v_p(2)$.
We treat both the proper-fraction and polynomial-quotient cases.

First suppose $l\le N$. Direct multiplication at $-i^2$ gives
\begin{equation}\label{eq:residue-derivative}
 D_N'(-i^2)=\frac{(-1)^{i-1}(N-i)!(N+i)!}{2i^2}.
\end{equation}
The factorial ratio appearing in $(N!)^2/D_N'(-i^2)$ is
\[
 \frac{(N!)^2}{(N-i)!(N+i)!}
 =\frac{\binom{2N}{N-i}}{\binom{2N}{N}}.
\]
For $0\le b\le a$, Legendre's factorial formula gives
\[
 v_p\binom ab
 =\sum_{j\ge1}\left(
 \floor{a/p^j}-\floor{b/p^j}-\floor{(a-b)/p^j}\right)
 \le\floor{\log_p a},
\]
since each summand is $0$ or $1$. Thus the displayed ratio has valuation
at least $-L_p$. The values $P_l(i^2)$ lose at most $v_p(2)$ by
\eqref{eq:central-binomial}; each $k_i(X)$ loses at most
$5L_p+2v_p(2)$, directly from \eqref{eq:formal-functional}.
The factor $2i^2$ in \eqref{eq:residue-derivative} causes no further
loss. These bounds control every residue. For $l<N$ there is no
polynomial quotient. For $l=N$ the sole quotient is constant and
contributes
\[
 (-1)^N\frac{\mu_0}{\binom{2N}{N}},\qquad \mu_0=\frac5{12}.
\]
All these contributions, including the product-basis loss, satisfy
\eqref{eq:small-prime-entry}.

Now suppose $l\ge N+1$. The denominator $D_N$ divides $P_l(-t)$
exactly. Up to a sign, the polynomial to which the right side of
\eqref{eq:bernoulli-moments} is applied is
\[
 \frac{(N!)^2}{(2l)!}\,x^7
 \prod_{a=N+1}^{l-1}(x-a)(x+a).
\]
Put $k=l-N-1\ge0$. At any integer $x$, each of the two products of
$k$ consecutive integers is divisible by $k!$. The remaining scalar
has the exact expression
\[
 \frac{(N!)^2(k!)^2}{(2l)!}
 =\frac1{l^2\binom{2l}{l}\binom{l-1}{N}^2}.
\]
Its valuation is at least $-5L_p$: the factor $l^2$ costs at most
$2L_p$, and the three binomial factors counted with multiplicity cost
at most $3L_p$. The polynomial degree is at most
$2l-2N+5\le4R-2N+1$. Equations
\eqref{eq:bernoulli-newton}--\eqref{eq:newton-derivative} show that
three derivatives followed by Bernoulli evaluation lose at most
$4L_p$ further powers, and division by $24$ costs $v_p(24)$.
Adding the product-basis loss $2v_p(2)$ proves the stated bound, since
$24\cdot2^2=96$.

Finally,
$\log E_N\le\log96+9\pi(N^{9/10})\log(4R+4)
=O(N^{9/10}\log N)$, using only the elementary bound $\pi(y)\le y$.
Since $r=6N/5$, this gives
$\log(E_N^r)=O(N^{19/10}\log N)=o(N^2)$.
\end{proof}

% END SHARED CORE 05_integer_valued

% BEGIN SHARED CORE 06_local_kernel
\section{The local quotient and restriction-kernel theorem}\label{sec:local-kernel}

The main local estimate must accommodate $R>N$: the entries then have
nonconstant polynomial quotients as well as proper fractions. We control
both contributions with the same restriction map.

Throughout this section a \emph{good prime} means a prime satisfying
\begin{equation}\label{eq:good-prime}
 p>5,\qquad p^2>4R+20.
\end{equation}
Write
\begin{equation}\label{eq:local-parameters}
 N=qp+\sigma,\quad 0\le\sigma<p,\qquad
 \epsilon=\floor{2\sigma/p},\qquad
 h=2q+1+\epsilon,\qquad D=\floor{(2R-2)/p}.
\end{equation}
In particular, $q=\floor{N/p}$ and $\epsilon\in\{0,1\}$.
Let $V$ be the free $\mathcal O_p$-module spanned by
$P_j(x^2)$, $0\le j<R$. Its elements are even polynomials, are
integer-valued at $p$, and have degree at most $2R-2<p^2$.

\begin{definition}\label{def:restriction-map}
For $1\le c\le(p-1)/2$, the \emph{reflected residue class} $c$ comprises
the physical indices $i\in\{1,\ldots,N\}$ with $i\equiv c$ or $-c$
modulo $p$. Equivalently, use the signed levels
\[
 K_c=\{a\in\Z:1\le |c+pa|\le N\},\qquad
 h_c=|K_c|=2q+\mathbf1_{c\le\sigma}+\mathbf1_{p-c\le\sigma}.
\]
Here $\mathbf1_E$ is $1$ when $E$ holds and $0$ otherwise.
Call a class \emph{maximal} when $h_c=h$, and denote the set of these
classes by $\mathcal C_p$. Their number is
\begin{equation}\label{eq:maximal-class-count}
 w_p=|\mathcal C_p|=
 \begin{cases}
 \sigma,&\epsilon=0,\\
 \sigma-(p-1)/2,&\epsilon=1.
 \end{cases}
\end{equation}
The set may be empty; when $\sigma=0$, the value $h$ is an upper
threshold, not an attained class size. Define the polynomial restriction
map
\begin{equation}\label{eq:restriction-map}
 \mathcal R_p:V/pV\longrightarrow
       \bigoplus_{c\in\mathcal C_p}\Fp[Y]_{\le D},\qquad
 \overline f\longmapsto
       \bigl(\overline{f(c+pY)}\bigr)_{c\in\mathcal C_p}.
\end{equation}
Reduction here is reduction of polynomial coefficients, not just
evaluation as a function on the finite field.
\end{definition}

The map is well-defined by \cref{lem:newton-tools}(iii). In particular,
$\rank\mathcal R_p\le(D+1)w_p$. Formula
\eqref{eq:maximal-class-count} follows by counting $c\le\sigma$ when
$\sigma<p/2$, and counting $p-\sigma\le c\le(p-1)/2$ when
$\sigma>p/2$.

\begin{proposition}\label{prop:local-kernel}
For every good prime, $p^{5+\epsilon}A_{N,R}(X)$ belongs to
$M_R(\mathcal O_p[X])$. Its reduction modulo $p$ is constant in $X$,
and, as a bilinear form on $V/pV$, has the kernel of $\mathcal R_p$
in its radical. In particular its rank is at most $(D+1)w_p$, and
\begin{equation}\label{eq:local-determinant-bound}
 v_p(F_N)\ge-(4+\epsilon)r-\min\{r,(D+1)w_p\}.
\end{equation}
The radical of a bilinear form consists of those vectors whose pairing
with every vector is zero.
\end{proposition}
\begin{proof}
For $f=\sum_j a_jP_j(x^2)\in V$, write
$\widehat f(t)=\sum_j a_jB_j(t)$, so $f(x)=\widehat f(-x^2)$.
The bilinear form represented by the full matrix is
$(N!)^2\J_X(\widehat f\widehat g/D_N)$.

\emph{The exact quotient.}
Expand $fg$ in the $P_l(x^2)$ basis, with coefficients in
$\mathcal O_p$ by \cref{lem:newton-tools}, and retain the terms
$0\le l\le N$; call this truncation $I$. It is even,
integer-valued at $p$, and has degree at most $2N$.
For each $i=0,1,\ldots,N$, all omitted terms vanish at $i$.
Thus $fg-I$ vanishes at $0,\pm1,\ldots,\pm N$.
Evenness gives a double zero at zero, so $x^2D_N(-x^2)$ divides
$fg-I$ in $\Q[x]$. In particular,
\begin{equation}\label{eq:normalized-local-quotient}
 Q_{f,g}(x)=(N!)^2\frac{f(x)g(x)-I(x)}{D_N(-x^2)}
\end{equation}
is an even polynomial. The contribution of this polynomial to the
bilinear form is obtained from \eqref{eq:bernoulli-moments}, with
$q(-x^2)=Q_{f,g}(x)$. The only polynomial quotient contributed by
$I$ is the constant from its $P_N$ term. Its valuation is at least
$-\epsilon$, because
$v_p((N!)^2/(2N)!)=-\epsilon$ and $\mu_0=5/12$ is a $p$-unit.

\emph{Division on each residue class.}
For nonzero $c$, factor exactly
\begin{equation}\label{eq:local-denominator-factorization}
 D_N(-(c+pY)^2)=p^{h_c}L_c(Y)U_c(Y),\qquad
 L_c(Y)=\prod_{a\in K_c}(Y-a),
\end{equation}
where $U_c\in\mathcal O_p[Y]$ has nonzero constant reduction.
Indeed, each physical index in the class supplies exactly one linear
factor $p(Y-a)$ up to sign; every remaining factor reduces to a
nonzero constant. The sign is included in $U_c$.
The substituted numerator $(fg-I)(c+pY)$ belongs to
$\mathcal O_p[Y]$. It is divisible by the monic polynomial $L_c$,
so its quotient by $L_c$ still belongs to $\mathcal O_p[Y]$.
This follows directly from the usual monic polynomial-division
algorithm. The remaining quotient by $U_c$ is an actual polynomial,
because \eqref{eq:normalized-local-quotient} already exists in
$\Q[x]$. The Gauss identity \eqref{eq:gauss-laws} preserves its lower
valuation bound, since $v_p(U_c)=0$. Finally
$v_p((N!)^2)=2q$. Hence
\[
 v_p(Q_{f,g}(c+pY))\ge2q-h_c.
\]
At residue zero the denominator instead has the form
$p^{2q}\prod_{a=1}^q(Y^2-a^2)U_0(Y)$, up to a sign, with
$\overline U_0$ a nonzero constant. The same monic-division argument
gives $v_p(Q_{f,g}(pY))\ge0$. Evenness covers the negative residue
classes. It follows that $p^{1+\epsilon}Q_{f,g}$ is integer-valued
at $p$.

The degree of $x^5Q_{f,g}$ is at most $4R-2N+1$.
This degree plus one is strictly less than $p^2$ by
\eqref{eq:good-prime}. In the Newton basis, each of the three
derivatives and the Bernoulli evaluation in
\eqref{eq:bernoulli-moments} therefore loses at most one power of $p$.
The denominator $24$ is a unit. Consequently the quotient contribution
has valuation at least $-(5+\epsilon)$.

\emph{The extra power on the restriction kernel.}
Suppose $\overline f\in\ker\mathcal R_p$. On a maximal class,
$\overline{fg(c+pY)}=0$. The reduction of $I(c+pY)$ has degree at most
$\floor{2N/p}=h-1$ by \cref{lem:newton-tools}(iii), and vanishes at
each of the $h$ levels in $K_c$, because $I$ agrees with $fg$ there.
These levels are distinct modulo $p$: their pairwise differences have
absolute value at most $2N/p<p$, by \eqref{eq:good-prime}.
The reduction of $I(c+pY)$ is therefore the zero polynomial.
The numerator in \eqref{eq:normalized-local-quotient} gains a factor
$p$ on every maximal class. Monic division and the subsequent
Gauss-valuation argument both preserve that factor. Nonmaximal
classes already satisfy $h_c\le h-1$.
Thus $p^\epsilon Q_{f,g}$ is integer-valued at $p$, and its moment
contribution has valuation at least $-(4+\epsilon)$.

\emph{The proper fractions.}
At a physical node in nonzero class $c$, every other node in that
class contributes exactly one power of $p$ to $D_N'(-i^2)$; all
other factors are units. Exactness follows also from $2N<p^2$,
which excludes a second power in a nonzero difference $i\pm j$.
Hence
\[
 v_p\left(\frac{(N!)^2}{D_N'(-i^2)}\right)=2q-h_c+1.
\]
The factor $k_i(X)$ has valuation at least $-5$, because all integers
in its harmonic sum are less than $p^2$. Only maximal classes can
reach valuation $-(5+\epsilon)$. On those classes the condition
$\overline f\in\ker\mathcal R_p$ gives $f(i)\in p\mathcal O_p$
and therefore the same extra power as for the quotient.
For a node $i=ap$, $1\le a\le q$, the derivative denominator has
valuation $2(q-1)$, the normalized residue factor has valuation $2$,
and $k_i$ has valuation at least $-1$. Such contributions are integral.
Finally, the coefficient of $X$ in $k_i$ is $i^4$; the preceding
formulas show that the coefficient of $X$ in the whole bilinear form
has valuation at least $-\epsilon$.

We have proved coefficientwise integrality after multiplication by
$p^{5+\epsilon}$, independence of $X$ in its reduction, and the
radical assertion. Rank is at most the dimension of the restriction
image. The same upper bound holds after restricting to the window.
If its reduced rank is $k\le\min\{r,(D+1)w_p\}$, constant invertible
row operations over $\mathcal O_p$ make $r-k$ rows zero modulo $p$.
Because the reduction is independent of $X$, every coefficient in
each such row is divisible by $p$. The determinant of the scaled
matrix therefore has valuation at least $r-k$. Removing the scale
$p^{5+\epsilon}$ from its $r$ rows proves
\eqref{eq:local-determinant-bound}.
\end{proof}

% END SHARED CORE 06_local_kernel

% BEGIN SHARED CORE 07_monomial
\section{Monomial residues, quotient support, and mixed minors}\label{sec:monomial}

We now use the matrix $H$ in \cref{lem:basis-change}. Write
$t_i=-i^2$ and $v_i=(1,t_i,\ldots,t_i^{r-1})^{\mathsf T}$.
Polynomial division followed by partial fractions gives the exact
decomposition
\begin{equation}\label{eq:monomial-residue-decomposition}
 H(X)=Q+\sum_{i=1}^N w_i(X)v_iv_i^{\mathsf T},\qquad
 w_i(X)=\frac{C_m(t_i)^2k_i(X)}{D_N'(t_i)}.
\end{equation}
The matrix $Q$ consists of the moments of the polynomial quotients.
It is constant in $X$. A summand $w_iv_iv_i^{\mathsf T}$ will be
called a \emph{residue atom}; this denotes precisely this rank-one
matrix, including its scalar weight.

\begin{lemma}\label{lem:monomial-support}
For a good prime, the moments needed in $Q$ have valuation at least
$-1$. They are integral whenever their index $k$ is less than $2p-3$.
There is a decomposition
$Q=Q_0+p^{-1}Q_1$ with $Q_0,Q_1\in M_r(\mathcal O_p)$ such that
$Q_1$ is supported in its last $t_p$ rows and columns, where
\begin{equation}\label{eq:quotient-terminal-size}
 t_p=\min\{r,\pos{2R-N-2p+2}\}.
\end{equation}
Consequently every minor of $Q$ of order $a$ has valuation at least
$-\min(a,t_p)$.
\end{lemma}
\begin{proof}
Whenever $n+1<p^2$, the Newton formula
\eqref{eq:bernoulli-newton}, applied to the integer-valued monomial
$x^n$, gives $v_p(\mathrm B_n)\ge-1$.
For $n\ge1$ in this range, \eqref{eq:bernoulli-distribution} gives
\[
 p\mathrm B_n\equiv\sum_{c=0}^{p-1}c^n\pmod p.
\]
Indeed, in expanding $(c+pY)^n$, the term of degree one in $Y$
has a factor $p\mathrm B_1=-p/2$, and every term of higher degree
$a$ has valuation at least $a-1\ge1$.
The sum of the $n$th powers in $\Fp^\times$ is zero unless
$(p-1)\mid n$, in which case it is $-1$. This follows by summing
the geometric progression in a generator of the cyclic group
$\Fp^\times$. Hence $\mathrm B_n$ is integral unless $(p-1)\mid n$.

For $n=2k+2$, the first three possible poles occur at
$n=p-1,2(p-1),3(p-1)$. In \eqref{eq:rational-moments} they are
canceled respectively by the factors $2k+3=p$, $2k+4=2p$, and
$2k+5=3p$. Thus the first uncanceled possible index is
$k=2p-3$. Every needed index satisfies $n+1<p^2$, because the
maximum quotient degree is $2R-N-2$ and
$p^2>4R+20$. The factor $24$ is a unit. This proves the moment
assertions with the required index precision.

The numerator $C_m(t)^2t^{k+l}$ and the monic denominator $D_N(t)$
have integer coefficients. Their polynomial quotient therefore has
integer coefficients and degree at most $2m+k+l-N$; if this upper
bound is negative the quotient is zero. A fractional entry can occur
only if
\[
 k+l\ge N+2p-3-2m.
\]
Since $k,l\le r-1$, this forces both $k$ and $l$ to be at least
$r-t_p$. This is exactly the terminal support asserted in the lemma.
Set $Q_0$ equal to $Q$ outside that terminal block and zero inside;
put $Q_1=p(Q-Q_0)$. These are integral matrices of the stated form.
In the multilinear column expansion of a minor of order $a$, at
most $\min(a,t_p)$ columns can contribute from $Q_1$ to a nonzero
term, because $\rank Q_1\le t_p$. Its denominator is accordingly
bounded by $p^{\min(a,t_p)}$.
\end{proof}

\begin{lemma}\label{lem:mixed-minors}
In the determinant expansion of \eqref{eq:monomial-residue-decomposition},
a term using $\ell$ distinct atoms is a product, up to sign, of
their $\ell$ weights, two evaluation minors of order $\ell$, and
a complementary quotient minor of order $r-\ell$.
Each evaluation minor is divisible in $\Z$ by the Vandermonde
product of its selected nodes.

Suppose $p\le N$ is good. In a nonzero reflected class containing
two nodes below $p$ and $h$ nodes above $p$, selecting $\ell$
atoms of which $b$ are above $p$ has denominator cost at most
\begin{equation}\label{eq:cluster-cost}
 5b+\ell(h+2-\ell)
 \le c_h(\ell):=5\min(\ell,h)+\ell(h+2-\ell),
 \qquad 0\le\ell\le h+2.
\end{equation}
Here a cost $e$ means a valuation lower bound $-e$, and negative
costs are permitted. The quotient cost in that same determinant
term is at most $\min(t_p,r-\ell_{\rm all})$, where
$\ell_{\rm all}$ is the total number of atoms selected from all classes.
\end{lemma}
\begin{proof}
Expand the determinant by multilinearity, or apply Cauchy--Binet and
then Laplace expansion to $Q+V\operatorname{diag}(w_i)V^{\mathsf T}$,
where $V$ has columns $v_i$. A repeated atom gives two proportional
columns and contributes zero. For each set $I$ of $\ell$ distinct
atoms and sets $J,K$ of $\ell$ matrix indices, the resulting term
has the form
\[
 \pm\left(\prod_{i\in I}w_i\right)
 \det V_{J,I}\det V_{K,I}\det Q_{J^c,K^c}.
\]
An evaluation minor is an alternating polynomial with integer
coefficients in its node variables. Setting any two variables equal
makes it zero, so it is divisible by their difference. The distinct
linear differences are prime, nonassociate factors in the polynomial
ring over $\Z$; their product therefore divides the minor over $\Z$.
Evaluation at the integer nodes proves the stated divisibility.

Each nonzero reflected class contains the two low nodes $c,p-c$.
Call its remaining $h$ nodes \emph{high}. Every difference of two
squared nodes in this class has valuation exactly one, since its
linear factors $i-j,i+j$ have absolute value below $p^2$, and only
one of them is divisible by $p$. Thus $D_N'(t_i)$ has valuation
$h+1$ for each node in this class. A low node has $v_p(k_i)\ge0$;
a high node has $v_p(k_i)\ge-5$. Ignoring, for now, the integral
factor $C_m(t_i)^2$, the product of weights costs at most
$5b+\ell(h+1)$. The two evaluation minors restore
$2\binom\ell2=\ell(\ell-1)$ powers of $p$. Subtraction gives
\eqref{eq:cluster-cost}. The complementary minor has order
$r-\ell_{\rm all}$, so \cref{lem:monomial-support} gives its stated
cost. In particular this cost and the atomic costs share the same
matrix dimension; they cannot be maximized independently.
\end{proof}

\begin{lemma}\label{lem:large-prime-integrality}
For a good prime, an all-low reflected pair contributes an integral matrix over
$\mathcal O_p[X]$. For every prime $p>N$, the entire monomial matrix
$H(X)$ is integral over $\mathcal O_p[X]$. Consequently the
denominators of $F_N$ involve only primes at most $2R$.
\end{lemma}
\begin{proof}
For $1\le i<p$ and $j=p-i$, the sum of inverse fifth powers of
the nonzero residues is zero modulo $p$, since $p>5$. Hence
$H_{p-i-1}^{(5)}\equiv H_i^{(5)}\pmod p$, and direct substitution
in \eqref{eq:formal-functional} gives
\begin{equation}\label{eq:stieltjes-reflection}
 k_{p-i}(X)\equiv k_i(X)\pmod p.
\end{equation}
For clarity, the extra term $i^{-5}$ in $H_i^{(5)}$ contributes
$-i^{-1}$ after multiplication by $i^4$, while the reflected term
$-1/(2(p-i))$ contributes $+1/(2i)$. Their sum is $-1/(2i)$,
as required.

Suppose this pair is the whole reflected class. Write
$D_N(t)=(t-t_i)(t-t_j)E(t)$. Both $E(t_i)$ and $E(t_j)$ are units
and are congruent modulo $p$. For any matrix entry the paired sum is
\[
 \frac{C_m(t_i)^2t_i^{k+l}k_i(X)/E(t_i)
       -C_m(t_j)^2t_j^{k+l}k_j(X)/E(t_j)}{t_i-t_j}.
\]
The numerator is zero modulo $p$, by \eqref{eq:stieltjes-reflection}
and $t_i\equiv t_j$. The denominator has valuation exactly one,
so the expression is integral coefficientwise.

For $p>N$, every node is low, and every reflected class is either
a singleton or such a pair. Singletons have unit derivative
denominators and hence give integral atoms. There are no zero-class
nodes. The quotient degree is at most
$2R-N-2<2p-3$, so \cref{lem:monomial-support} makes the quotient
integral as well. These primes are good because $N\ge20$ and
$R=5N/4$. Finally, for $p>2R$ the factorial prefactor in
\eqref{eq:monomial-change} is also integral (indeed a unit).
This proves the denominator-support assertion.
\end{proof}

When integral pairs or other integral atoms are absorbed into the
quotient matrix, its integral part $Q_0$ may become affine in $X$.
This causes no change to the argument: we then have
$Q_0\in M_r(\mathcal O_p[X])$, while the fractional part $p^{-1}Q_1$
retains its terminal support. The minor estimate and all mixed-minor
valuations remain coefficientwise polynomial assertions.

% END SHARED CORE 07_monomial

% BEGIN SHARED CORE 08_root_gain
\section{The common-root gain and its shared dimension bound}\label{sec:root-gain}

The common factor $C_m$ improves the residue estimate in a small,
explicit collection of reflected classes. We prove the gain before
passing to any limiting prime distribution, so that its interaction
with the polynomial quotient is exact.

In this section assume
\begin{equation}\label{eq:root-gain-range}
 N/3<p\le N,\qquad p>2m,\qquad p>5,\qquad p^2>4R+20.
\end{equation}
Thus $q=\floor{N/p}$ is $1$ or $2$, and write again
$\sigma=N-qp$. Define
\begin{equation}\label{eq:high-node-count}
 H_p=N-p-(q-1),\qquad
 t_p=\min\{r,\pos{2R-N-2p+2}\},
\end{equation}
and
\begin{equation}\label{eq:affected-node-count}
 A_m(p)=2(q-1)(m-1)+\min(m-1,\sigma)+\pos{m-p+\sigma}.
\end{equation}
Here $H_p$ counts high physical nodes in nonzero reflected classes.
We call a class \emph{affected} if its representative satisfies
$1\le c<m$; $A_m(p)$ will count its high nodes.

\begin{proposition}\label{prop:root-gain}
Under \eqref{eq:root-gain-range}, and provided $r\ge H_p$,
\begin{equation}\label{eq:root-gain-bound}
 v_p(\det H(X))\ge
 -7H_p+2A_m(p)-\min(t_p,r-H_p).
\end{equation}
The condition $r\ge H_p$ holds for our fixed parameters.
All valuations in this statement are coefficientwise in $X$.
\end{proposition}
\begin{proof}
\emph{Valuation of the common roots.}
Consider a nonzero reflected class represented by $c<m$.
The low atom at $i=c$ vanishes exactly, because $C_m(-c^2)=0$.
At every other node $i$ of the class, exactly one factor in
\[
 C_m(-i^2)=\prod_{a=0}^{m-1}(i^2-a^2)
\]
is divisible by $p$, namely the factor with $a=c$.
Indeed, $p>2m$ makes the residues $\pm a$, $0\le a<m$,
distinct except at zero. The surviving divisible factor has valuation
exactly one: the nonzero integer $i-c$ or $i+c$ that contains $p$
has absolute value below $p^2$. The other linear factor is a unit.
Thus $C_m(-i^2)^2$ contributes exactly two powers of $p$ at each
surviving node of an affected class.

Write $h$ for the number of high nodes in the particular class under
discussion, so that the original class has $2+h$ nodes.
For $q=1$ its possible values are $0,1,2$; for $q=2$ they are
$2,3,4$. A mixed-minor term selecting $\ell$ atoms from an affected
class, of which $b$ are high, therefore has cost at most
\begin{equation}\label{eq:affected-cluster-cost}
 5b+\ell(h+1)-2\ell-\ell(\ell-1)
 =5b+\ell(h-\ell).
\end{equation}
There is at most one surviving low atom, so $0\le\ell\le h+1$.
At fixed $\ell$ this is bounded above by
\[
 c_h^{\rm aff}(\ell)=5\min(\ell,h)+\ell(h-\ell).
\]
For $1\le h\le4$, its maximum is $5h$, attained at $\ell=h$.
If $\ell=h-d<h$, its deficit from that maximum is
$d(5-h)+d^2\ge d$. If $\ell=h+1$, its value is smaller than
$5h$. In an unaffected class, \eqref{eq:cluster-cost} instead gives
maximum $7h$, and the deficit at $\ell=h-d<h$ is
$d(7-h)+d^2\ge d$. The choices $\ell=h+1,h+2$ do not exceed
$7h$ when $h\ge1$.

Classes with $h=0$ may be absorbed into the integral part of the
quotient matrix. An unaffected all-low pair is integral by
\cref{lem:large-prime-integrality}. In an affected class the node
$c$ vanishes, and the other low node has valuation at least one
after the two-power root gain. Absorption preserves the terminal
fractional support, as explained at the end of \cref{sec:monomial}.

\emph{The zero class.}
For $i=ap$, $1\le a\le q$, only the factor with root zero in
$C_m(-i^2)$ is divisible by $p$; it has valuation two.
Thus $v_p(C_m(-i^2)^2)=4$. The derivative denominator has valuation
$2(q-1)$, while $v_p(k_i(X))\ge-1$. The corresponding atom has
valuation at least $5-2q\ge1$. All zero-class atoms can therefore
be absorbed individually into the integral part of the quotient.

\emph{Exact counts.}
In a class $c$, the number of high nodes is
\[
 h(c)=2(q-1)+\mathbf1_{c\le\sigma}+\mathbf1_{c\ge p-\sigma}.
\]
Summing this identity over $1\le c<m$ gives exactly
\eqref{eq:affected-node-count}. There are $N-p$ physical nodes
strictly above $p$, of which $q-1$ are in the zero class. This proves
\eqref{eq:high-node-count}. The sum of the maximum class costs is
therefore
\[
 M_p=7\bigl(H_p-A_m(p)\bigr)+5A_m(p)
     =7H_p-2A_m(p).
\]

\emph{The shared dimension.}
After the integral classes and atoms have been absorbed, consider
one mixed-minor term with $\ell$ atoms in total. The preceding
deficit bounds imply that its atomic cost is at most
\[
 M_p-\pos{H_p-\ell}.
\]
Indeed, summing the deficits $(h(c)-\ell_c)_+$ over the classes
is at least $(\sum_ch(c)-\sum_c\ell_c)_+$.
The remaining minor has size $r-\ell$, so its quotient cost is
at most $\min(t_p,r-\ell)$. For $r\ge H_p$,
\[
 \min(t_p,r-\ell)-\pos{H_p-\ell}
 \le\min(t_p,r-H_p).
\]
If $\ell\ge H_p$, this follows by monotonicity. If $\ell<H_p$,
write $d=H_p-\ell$ and use
$\min(t_p,r-H_p+d)-d\le\min(t_p,r-H_p)$.
Every term thus has cost at most $M_p+\min(t_p,r-H_p)$.
The nonarchimedean inequality for Gauss valuations preserves this
bound upon summation and proves \eqref{eq:root-gain-bound}.
\end{proof}

\subsection{Restoring the factorial basis}
For a good prime, the exact opposite valuation of the prefactor in
\eqref{eq:monomial-change} is
\begin{equation}\label{eq:window-factorial-cost}
 f_p^{[m,R)}=2\sum_{j=m}^{R-1}\floor{2j/p}-2rq.
\end{equation}
Here there are no higher factorial terms because $2R<p^2$.
Combining this with \cref{prop:root-gain} gives the signed clearing
exponent
\begin{equation}\label{eq:upper-band-exponent}
 e_p=f_p^{[m,R)}+7H_p-2A_m(p)+\min(t_p,r-H_p),
 \qquad v_p(F_N)\ge-e_p.
\end{equation}
Using a negative $e_p$ is legitimate: it records a common factor
of all coefficients, as in \cref{lem:primitive-normalizer}.

\subsection{The size of the root gain on the upper prime band}
For $N/2<p<N$, put $\beta=p/N$. Then $q=1$, and the exact count
\eqref{eq:affected-node-count} gives, uniformly on this interval,
\begin{equation}\label{eq:affected-mass}
 \frac{A_m(p)}N
 =\min(s,1-\beta)+\pos{1+s-2\beta}+O(1/N).
\end{equation}
Relative to the bound that discards the root gain, the limiting
clearing-exponent profile is reduced by twice this expression.
Its integral reduction is the exact quantity
\begin{equation}\label{eq:integrated-root-gain}
 2\int_{1/2}^1
 \left[\min(s,1-\beta)+\pos{1+s-2\beta}\right]\dd\beta
 =s-\frac{s^2}{2}=\frac{39}{800}.
\end{equation}
To check the integral, the first summand has integral
$s/2-s^2/2$, and the second is a triangle of base $s/2$ and
height $s$, hence has integral $s^2/4$.

The proposition also proves a gain on $N/3<p<N/2$, but our final
arithmetic estimate uses the weaker bound obtained by discarding
$A_m(p)$ on that band. The local floor formulas and their limiting
affine pieces are used on open prime-ratio intervals. Their boundary
treatment, including the possible jumps at rational endpoints, is
given explicitly in \cref{sec:arithmetic}.

% END SHARED CORE 08_root_gain

% BEGIN SHARED CORE 09_potential
\section{An explicit positive logarithmic potential}\label{sec:potential}

We now estimate the determinant at its positive integral value. This part
uses only elementary potential identities and a quantified construction of
trial polynomials, proved in the next section. In particular, no general
asymptotic theorem for orthogonal polynomials is needed.

\subsection{Exact factorization into monic norms}
Let $\widehat h_{k,N,m}$ be the least squared norm of a real monic polynomial
of degree $k$ for the measure
\[
 \frac{B_m(t)^2}{D_N(t)}\,\dd\mu(t).
\]
Here and below $\mu$ is understood in the variable $t=x^2$, as defined in
\cref{sec:functional}. A monic polynomial has leading coefficient $1$.
Existence and uniqueness of the minimizer follow by minimizing a strictly
positive quadratic form in its $k$ lower coefficients. Subtracting its
orthogonal projections onto lower degrees gives the same minimizer.

\begin{lemma}\label{lem:norm-product}
The exact determinant identity is
\begin{equation}\label{eq:norm-product}
 F_N(\zeta(5))=
 \prod_{k=0}^{r-1}
 \frac{(N!)^2((2m)!)^2}{((2m+2k)!)^2}\,
 \widehat h_{k,N,m}.
\end{equation}
\end{lemma}
\begin{proof}
The polynomial $B_{m+k}/B_m$ has degree $k$ and leading coefficient
$(-1)^k(2m)!/(2m+2k)!$. A triangular change to monic orthogonal polynomials
has determinant equal to the product of these leading coefficients.
A Gram determinant changes by the square of that determinant.
Multiplying by the entry scale $(N!)^2$ in each of the $r$ rows gives
\eqref{eq:norm-product}. The common factor $B_m$ remains in the measure
throughout this computation.
\end{proof}

For $a\ge0$ and $y\ge0$, define
\begin{equation}\label{eq:external-field}
 I_a(y)=\int_0^a\log(y^2+t^2)\,\dd t,\qquad
 V_s(y)=2\pi y+I_1(y)-2I_s(y).
\end{equation}
The integrals at $y=0$ converge, with the convention $0\log0=0$.
In particular $V_s(0)=-2-4s\log s+4s$.
Monotonicity of $t\mapsto\log(y^2+t^2)$ gives the exact Riemann-sum bounds
\begin{equation}\label{eq:riemann-weight}
 D_N(N^2y^2)\ge N^{2N}e^{NI_1(y)},\qquad
 |B_m(N^2y^2)|^2
 \le\frac{N^{4m}}{((2m)!)^2}e^{2NI_s(y)}.
\end{equation}
The second inequality at $y=0$ follows directly from $B_m(0)=0$.

\subsection{The measure and its mass}
For $b,t>0$, define a symmetric positive measure $\nu_{b,t}$ on $[-b,b]$
by its density
\begin{equation}\label{eq:nu-density}
 \frac{\dd\nu_{b,t}}{\dd x}
 =\frac{t\sqrt{b^2-x^2}}
 {2\pi\sqrt{t^2+b^2}(x^2+t^2)}
 \quad (|x|<b).
\end{equation}
Also define $\lambda_b$ by the density
$\pi^{-1}\arcosh(b/|x|)$ on $0<|x|<b$, where
$\arcosh z=\log(z+\sqrt{z^2-1})$ for $z\ge1$.
The densities have integrable singularities where indicated.

\begin{lemma}\label{lem:positive-potential}
The measure
\begin{equation}\label{eq:potential-measure}
 \eta_{b,s}
 =\lambda_b+\int_0^1\nu_{b,t}\,\dd t
                  -2\int_0^s\nu_{b,t}\,\dd t
 =2\int_s^\infty\nu_{b,t}\,\dd t+\int_0^1\nu_{b,t}\,\dd t
\end{equation}
is positive. Its total mass is
\begin{equation}\label{eq:u-of-b}
 u_s(b)=
 \frac{b+1-2s-\sqrt{1+b^2}+2\sqrt{s^2+b^2}}2.
\end{equation}
The function $u_s$ increases continuously and strictly from zero to
infinity. We denote its inverse by $b_s(u)$ for $u>0$.
\end{lemma}
\begin{proof}
Direct integration of \eqref{eq:nu-density} gives mass
$(1-t/\sqrt{t^2+b^2})/2$. For completeness, its Cauchy transform is
\begin{equation}\label{eq:cauchy-nu}
 \int_{-b}^{b}\frac{\dd\nu_{b,t}(x)}{z-x}
 =\frac{z-t\sqrt{z^2-b^2}/\sqrt{t^2+b^2}}{2(z^2+t^2)},
 \qquad z\notin[-b,b],
\end{equation}
where the branch is fixed by $\sqrt{z^2-b^2}\sim z$ at infinity.
The jump of the right side across the interval is $-2\pi i$ times
\eqref{eq:nu-density}; its apparent poles at $\pm it$ cancel. The difference
between the two sides extends analytically across the interval and its
endpoints and vanishes at infinity, so it is zero. At the endpoints this
removability follows also from the $O(\sqrt{b-|x|})$ density bound, which
makes the transforms bounded there. The coefficient of $1/z$
then gives the mass just stated. Alternatively, substitution $x=b\sin\theta$
gives that mass by a real integral.

At each $0<|x|<b$, substitution $z=\sqrt{t^2+b^2}$ in the density integral
gives $\lambda_b=2\int_0^\infty\nu_{b,t}\,\dd t$. Tonelli's theorem applies
to these nonnegative densities. This proves the second expression in
\eqref{eq:potential-measure} and hence positivity. Integrating the masses
gives \eqref{eq:u-of-b}. Finally,
\[
 u_s'(b)=\frac12\left(
  1-\frac b{\sqrt{1+b^2}}+\frac{2b}{\sqrt{s^2+b^2}}
 \right)>0.
\]
The formula shows $u_s(0)=0$ and $u_s(b)\to\infty$.
\end{proof}

We use $\eta_{b,s}$, rather than $\mu$, for this auxiliary measure to avoid
confusing it with the original measure in the determinant. Define its
logarithmic potential by
\[
 U_{b,s}(y)=\int_{-b}^{b}\log|y-x|\,\dd\eta_{b,s}(x).
\]
The logarithmic singularity is integrable against its density; the potential
is continuous on the real line. To justify the latter statement explicitly,
$0\le\eta_{b,s}\le(3/2)\lambda_b$, whose density is
$O(1+|\log|x||)$ for fixed $b$. Integrating the logarithmic kernel over a
short interval around its singularity then gives a bound tending to zero
uniformly in the center of that interval; off that interval continuity
follows by dominated convergence.

\begin{proposition}\label{prop:potential-inequality}
For $u=u_s(b)$ there is a constant $\ell_s(u)$ such that
\begin{align}
 4U_{b,s}(y)-V_s(y)&=\ell_s(u) &&(0\le y\le b),\label{eq:potential-equality}\\
 4U_{b,s}(y)-V_s(y)&\le\ell_s(u) &&(y\ge0).\label{eq:potential-inequality}
\end{align}
Moreover,
\begin{align}
 \ell_s(u)={}&2+4s\log s-4s
  +4u\bigl(\log(b/2)-1\bigr)\notag\\
 &-2\log\frac{1+\sqrt{1+b^2}}2
  +4s\log\frac{s+\sqrt{s^2+b^2}}{2s},
 \qquad b=b_s(u).\label{eq:ell-explicit}
\end{align}
\end{proposition}
\begin{proof}
On $0<y<b$, the real part of \eqref{eq:cauchy-nu} is
$y/(2(y^2+t^2))$. Its principal-value meaning is obtained by differentiating
the logarithmic potential; the density is smooth locally away from zero
and the endpoints. Using \eqref{eq:potential-measure} therefore gives
$4U_{b,s}'(y)=V_s'(y)$. For $y>b$, the same identity gives
\begin{align}
 4U_{b,s}'(y)-V_s'(y)
 ={}&-4\int_s^\infty
 \frac{t\sqrt{y^2-b^2}}{\sqrt{t^2+b^2}(y^2+t^2)}\,\dd t\notag\\
 &-2\int_0^1
 \frac{t\sqrt{y^2-b^2}}{\sqrt{t^2+b^2}(y^2+t^2)}\,\dd t<0.
 \label{eq:tail-derivative}
\end{align}
For example, the part without square roots in the first line before
subtraction is $4\int_s^\infty y/(y^2+t^2)\,\dd t
=2\pi-4\arctan(s/y)$; the other part is $2\arctan(1/y)$,
exactly $V_s'(y)$. Continuity at $0$ and $b$ extends the identities to
their endpoints and proves \eqref{eq:potential-equality}--\eqref{eq:potential-inequality}.

Let $\omega_b$ be the unit arcsine measure with density
$1/(\pi\sqrt{b^2-x^2})$ on $(-b,b)$. Differentiating
\eqref{eq:nu-density} with respect to $b$ gives its mass derivative times
$\omega_b$; differentiating the density of $\lambda_b$ gives $\omega_b$.
The boundary density is zero, so there is no boundary atom. It follows
against bounded continuous test functions that
\[
 \frac{\partial\eta_{b,s}}{\partial b}=u_s'(b)\omega_b.
\]
The arcsine potential on its interval is $\log(b/2)$: this follows from
$x=b\cos\theta$ and
$\pi^{-1}\int_0^\pi\log|y-b\cos\theta|\,\dd\theta=\log(b/2)$
for $|y|\le b$. For the kernel $\log|x|$, first use bounded truncations.
The explicit density derivatives have locally uniform finite logarithmic
moments for $b$ in a compact subinterval of $(0,\infty)$, so the truncations
may be removed by dominated convergence. Evaluating
\eqref{eq:potential-equality} at $y=0$ and differentiating therefore gives
\[
 \frac{\dd\ell_s}{\dd u}=4\log(b_s(u)/2),\qquad
 \ell_s(0)=-V_s(0).
\]
For the initial value, evaluate the potential at $y=0$.
The domination $0\le\eta_{b,s}\le(3/2)\lambda_b$, followed by $x=bz$, gives
\[
 |U_{b,s}(0)|\le Cb(1+|\log b|)\longrightarrow0
 \qquad(b\downarrow0);
\]
the required fixed integrals of $\arcosh(1/|z|)$ and its product with
$|\log|z||$ over $(-1,1)$ are finite. This proves the initial value
without relying on a small-mass argument alone.
Integration by parts reduces the constant to
$4u\log(b/2)-4\int_0^b u_s(v)/v\,\dd v$ plus the initial value. The elementary
identity
\[
 \int_0^b\frac{u_s(v)}v\,\dd v
 =u_s(b)+\frac12\log\frac{1+\sqrt{1+b^2}}2
       -s\log\frac{s+\sqrt{s^2+b^2}}{2s}
\]
is checked by differentiation and its value at $b=0$. Substitution proves
\eqref{eq:ell-explicit}.
\end{proof}

Uniformly for $u$ in a fixed compact subinterval of $(0,\delta]$, the
quantity on the left of \eqref{eq:tail-derivative} is bounded above by a
fixed negative constant outside a fixed compact interval. Indeed,
$U_{b,s}'(y)\le u/(y-b)$ there and
$V_s'(y)\ge2\pi-4s/y$. This uniform tail estimate will justify integrating
our pointwise polynomial bounds.

% END SHARED CORE 09_potential

% BEGIN SHARED CORE 10_discretization
\section{Quantified discretization and determinant decay}\label{sec:discretization}

The passage from a continuous potential to a polynomial must control the
logarithmic singularity uniformly. We provide that control explicitly.

\begin{lemma}\label{lem:quantile-bound}
Fix $0<\varepsilon<\delta$. For all sufficiently large $N$ and every integer $k$ with
$\varepsilon N\le k\le\delta N$, choose $b=b_s(k/N)$. There is a real monic
polynomial $p_{k,N}$ of degree $k$ such that, uniformly for $y\ge0$,
\begin{equation}\label{eq:trial-polynomial}
 |p_{k,N}(N^2y^2)|^2
 \le N^{4k+C_\varepsilon}\exp\{4NU_{b,s}(y)\}.
\end{equation}
\end{lemma}
\begin{proof}
Let $\sigma_N$ be the image of $\eta_{b,s}$ under $x\mapsto x^2$.
Thus $\sigma_N(E)=\eta_{b,s}(\{x:x^2\in E\})$ for a Borel set $E$;
this operation preserves its total mass $k/N$. Symmetry gives
\begin{equation}\label{eq:fold-potential}
 \int\log|y^2-v|\,\dd\sigma_N(v)=2U_{b,s}(y).
\end{equation}
On the present compact parameter range its density is bounded by
\begin{equation}\label{eq:squared-density}
 C v^{-1/2}(1+|\log v|),\qquad 0<v\le b^2.
\end{equation}
Indeed, \eqref{eq:potential-measure} is bounded above by
$\lambda_b+\int_0^1\nu_{b,t}\,\dd t$. The first density is
$O(1+|\log|x||)$, and the second is bounded by a constant times
$\int_0^1 t/(x^2+t^2)\,\dd t$, which has the same bound. The change of
variable $v=x^2$ proves \eqref{eq:squared-density}.

Partition $[0,b^2]$ into $k$ consecutive intervals, each of
$\sigma_N$-mass $1/N$, and take their ordinary coordinate midpoints $v_j$.
The measure has positive density in the interior and no atoms, so these
quantile intervals exist. From \eqref{eq:squared-density}, every interval
of length $a\le1$ has mass at most $C\sqrt a(1+|\log a|)$.
Consequently each quantile interval has length at least $N^{-3}$ for all
sufficiently large $N$, uniformly in $k$.

We claim the upper quadrature estimate
\begin{equation}\label{eq:log-quadrature}
 \sum_{j=1}^k\log|v-v_j|
 \le N\int\log|v-w|\,\dd\sigma_N(w)+C\log N
 \qquad(v\ge0).
\end{equation}
When $v$ is in the support, remove its containing interval and at most two
neighbors on each side. When $v$ is outside the support but in a fixed
bounded interval, remove the nearest terminal interval and its two inward
neighbors. On each remaining side $\log|v-w|$ is monotone. On one cell, its
sampled value exceeds its $N$-weighted integral by no more than the
oscillation on that cell. Summing these oscillations telescopes to the
total variation on that side. Distances to the remaining cells are at
least $N^{-3}$, so the total is $O(\log N)$.

For the union $E$ of the at most five removed cells, define
$\log^-z=\max(0,-\log z)$ for $z>0$. The interval-mass estimate and
integration of nonnegative level sets give
\[
 \int_E\log^-|v-w|\,\dd\sigma_N(w)
 \le\int_0^\infty
   \min\{5/N,\ C(1+t)e^{-t/2}\}\,\dd t
 =O(\log N/N).
\]
The positive parts of the sampled logarithms are bounded on a fixed
bounded interval. This proves \eqref{eq:log-quadrature} there. For $v$
beyond twice a fixed upper bound for the support, monotonicity gives
bounded total variation without removing cells. If a sample equals $v$,
the left side is $-\infty$ and the inequality is automatic.

Set $p_{k,N}(t)=\prod_{j=1}^k(t-N^2v_j)$. Taking twice
\eqref{eq:log-quadrature}, then using \eqref{eq:fold-potential}, gives
\eqref{eq:trial-polynomial}. The coefficients of this trial polynomial
need not be rational: it is used only to bound a minimum over real
monic polynomials.
\end{proof}

\begin{proposition}\label{prop:analytic-rate}
Define $b=b_s(u)$ for $u>0$ and
\begin{align}
 f_s(u)={}&
 4u\log\frac{b}{4(s+u)}
 +4s\log\frac{s+\sqrt{s^2+b^2}}{4(s+u)}
 -2\log\frac{1+\sqrt{1+b^2}}2,\label{eq:f-profile}\\
 f_s(0)={}&-4s\log2.\notag
\end{align}
Then $f_s$ is continuous on $[0,\delta]$, and
\begin{equation}\label{eq:analytic-rate}
 \limsup_{N\to\infty}\frac{\log F_N(\zeta(5))}{N^2}
 \le \int_0^\delta f_s(u)\,\dd u.
\end{equation}
\end{proposition}
\begin{proof}
The original density in the $x$ variable satisfies
\[
 0<\frac{\dd\mu(x^2)}{\dd x}
 =\frac{x^5}{12}\left(\frac{\dd}{\dd x}\right)^4
                 \frac1{e^{2\pi x}-1}
 \le C(1+x^5)e^{-2\pi x}.
\]
For $x\ge1$ this follows by differentiating the positive geometric series;
for $0<x\le1$ it follows from its Laurent expansion, since multiplication
by $x^5$ removes the leading pole. These arguments also show boundedness
at zero.

Combine this density bound, \eqref{eq:riemann-weight}, and
\eqref{eq:trial-polynomial}, substituting $x=Ny$. The exponent left in the
integral is $N(4U_{b,s}(y)-V_s(y))$.
By \cref{prop:potential-inequality} it is at most $N\ell_s(k/N)$.
Outside a fixed compact interval it decreases at a uniform linear rate,
by \eqref{eq:tail-derivative} and the following paragraph of
\cref{sec:potential}. All remaining factors are bounded by a fixed
power of $N$ times $1+y^5$. The variational definition of
$\widehat h_{k,N,m}$ therefore gives, uniformly when
$\varepsilon N\le k\le\delta N$,
\begin{equation}\label{eq:norm-bound}
 \log\widehat h_{k,N,m}
 \le(4k+4m-2N)\log N-2\log((2m)!)
       +N\ell_s(k/N)+O_\varepsilon(\log N).
\end{equation}
Using $\log(n!)=n\log n-n+O(\log(n+1))$ in
\eqref{eq:norm-product}, the logarithm of its $k$th factor is at most
\[
 N\{\ell_s(u)-2+4(s+u)-4(s+u)\log(2(s+u))\}
 +O_\varepsilon(\log N),\qquad u=k/N.
\]
The expression in braces equals \eqref{eq:f-profile} after substitution
of \eqref{eq:ell-explicit}. This calculation exhibits explicitly the
cancellation of all $N\log N$ terms and of the common $(2m)!$ factors.
Since $u_s(b)=b/2+O(b^2)$ near zero, its limit as $u\downarrow0$ is
$-4s\log2$.

It remains to treat $0\le k<\varepsilon N$, where the preceding uniformity
was not asserted. The actual polynomial $B_{m+k}/B_m$, after multiplication
by its inverse leading coefficient, is an admissible monic trial
polynomial. Hence the $k$th factor in \eqref{eq:norm-product} is at most
\[
 (N!)^2\int_0^\infty\frac{B_{m+k}(x^2)^2}{D_N(x^2)}
                         \,\dd\mu(x^2)
 \le\int_0^\infty B_{m+k}(x^2)^2\,\dd\mu(x^2).
\]
Here $D_N(x^2)\ge(N!)^2$. If $j=m+k$, then
$|B_j(x^2)|\le(x+j)^{2j}/(2j)!$. Substitution $z=x+j$, followed by extending
the integral down to zero, bounds the last display by
\[
 \frac{Ce^{2\pi j}}{((2j)!)^2}
 \left\{\frac{(4j)!}{(2\pi)^{4j+1}}
       +\frac{(4j+5)!}{(2\pi)^{4j+6}}\right\}
 \le e^{C_1N}N^{C_2}.
\]
The last inequality follows, for example, from
$(4j)!/((2j)!)^2=\binom{4j}{2j}\le2^{4j}$ and $j\le R$.
Thus these initial rows contribute at most
$C_1\varepsilon N^2+O(N\log N)$.

Summing the remaining rows and passing to a Riemann integral yields
$\limsup N^{-2}\log F_N(\zeta(5))
\le C_1\varepsilon+\int_\varepsilon^\delta f_s(u)\,\dd u$.
Finally let $\varepsilon\downarrow0$, using continuity at zero.
\end{proof}

% END SHARED CORE 10_discretization

% BEGIN SHARED CORE 11_arithmetic
\section{Summation of the local arithmetic bounds}\label{sec:arithmetic}

We turn the coefficientwise estimates into an upper bound for
$\log Q_N$. The \emph{arithmetic profile} below is a function of
$\beta=p/N$ giving the limiting denominator exponent divided by $N$.
Some of its values are negative: these record divisibility of every
coefficient and are retained in the primitive normalizer.

\subsection{The profile below 9/20}
For $0<\beta<9/20$, put
\[
 q=\floor{1/\beta},\qquad
 \theta=1/\beta-q,\qquad
 \epsilon=\floor{2\theta},\qquad
 D=\floor{2\rho/\beta}.
\]
The limit of \cref{prop:local-kernel} is
\begin{equation}\label{eq:rank-profile}
 e_{\mathrm{rank}}(\beta)
 =(4+\epsilon)\delta+
 \min\{\delta,(D+1)\beta(\theta-\epsilon/2)\}.
\end{equation}
Away from the floor discontinuities, its finite version differs from
$Ne_{\mathrm{rank}}(p/N)$ by $O(1)$ on any fixed compact $\beta$-interval.
The finite parameter in that proposition is
$\floor{(2R-2)/p}$; the possible difference from $D$ at its boundaries
will be handled explicitly below.

There is a convenient bound for the remaining interval near zero. Since
$D+1\le2\rho/\beta+1$ and $\theta-\epsilon/2\ge0$,
\[
 e_{\mathrm{rank}}(\beta)
 \le4\delta+2\rho\theta-s\epsilon
                   +\beta(\theta-\epsilon/2)
 \le4\delta+2\rho\theta+\frac{\beta}{2}.
\]
The interval $(0,1/40)$ is the union of the complete periods
$\beta=1/(q+\theta)$, $q\ge40$, $0<\theta<1$.
The weight $(q+\theta)^{-2}$ decreases with $\theta$, so its weighted
mean of $\theta$ is at most $1/2$. One proof pairs $\theta$ and $1-\theta$
in the integral of $(\theta-1/2)(q+\theta)^{-2}$.
Consequently
\begin{equation}\label{eq:small-profile-tail}
 \int_0^{1/40}e_{\mathrm{rank}}(\beta)\,\dd\beta
 \le\frac{4\delta+\rho}{40}
           +\frac{1}{4\cdot40^2}
 =\frac{969}{6400}.
\end{equation}

\subsection{An explicit table on the upper interval}
For $\beta>9/20$ the monomial estimates give a simpler table. The exact
factorial charge from the change of basis is
\[
 f_p=2\sum_{j=m}^{R-1}\floor{2j/p}
                  -2r\floor{N/p}.
\]
Since $p>2m$ here, its normalized limit is
\begin{equation}\label{eq:factorial-profile}
 f_{\mathrm{fac}}(\beta)
 =2D\rho-\frac{\beta D(D+1)}2-2\delta q,
 \qquad D=\floor{2\rho/\beta},\quad q=\floor{1/\beta}.
\end{equation}
For example,
$N^{-1}\sum_j\floor{2j/p}\to\int_s^\rho\floor{2x/\beta}\,\dd x$;
the integral equals
$D\rho-\beta D(D+1)/4$ because $s<\beta/2$.

For $9/20<\beta<1/2$, $q=2$, and the residue classes in
\cref{lem:mixed-minors} have $h=2,3$, or $4$ high atoms.
Their total number, divided by $N$, tends to $1-\beta$.
Apply \cref{prop:root-gain}, discarding its nonnegative term $2A_m(p)$;
that proposition also handles the zero class and the shared dimension.
The bound $7h$ for each class and the quotient support give the additional
cost
$7(1-\beta)+T(\beta)$, where
\[
 T(\beta)=(3/2-2\beta)_+.
\]
This uses the available matrix dimension, since
$T(\beta)\le\delta-(1-\beta)$ on the range in question.
For $1/2<\beta<1$, \cref{prop:root-gain} improves the cost to
\[
 7(1-\beta)-2A(\beta)+T(\beta),\qquad
 A(\beta)=\min(s,1-\beta)+(1+s-2\beta)_+ .
\]
The term $A$ is the limiting number of affected high atoms divided by $N$.
Substitution in \eqref{eq:factorial-profile} gives the following valid
profile on the \emph{open interiors} of the intervals:

\begin{center}
\begin{tabular}{@{}ll@{}}
\toprule
Interval for $\beta$ & $e_{\mathrm{upper}}(\beta)$\\
\midrule
$(9/20,\,1/2)$ & $81/5-24\beta$\\
$(1/2,\,21/40)$ & $139/10-15\beta$\\
$(21/40,\,5/8)$ & $16-19\beta$\\
$(5/8,\,3/4)$ & $27/2-15\beta$\\
$(3/4,\,5/6)$ & $12-13\beta$\\
$(5/6,\,19/20)$ & $19/2-10\beta$\\
$(19/20,\,1)$ & $38/5-8\beta$\\
\bottomrule
\end{tabular}
\end{center}

All entries follow from the displayed formulas; the changes occur precisely
where $D$, $T$, or $A$ changes formula. Their integral is
\begin{equation}\label{eq:upper-table-area}
 \int_{9/20}^{1}e_{\mathrm{upper}}(\beta)\,\dd\beta
 =\frac{901}{600}.
\end{equation}
In the $q=2$ interval we have discarded a nonnegative common-root
improvement, so no claim of optimality is needed.

For $p>N$, the monomial matrix is coefficientwise integral by
\cref{lem:large-prime-integrality}. Thus only the factorial charge remains.
For $x\le\rho=5/4<3/2$,
\[
 \int_1^\infty\floor{2x/\beta}\,\dd\beta
 =(2x-1)_++(x-1)_+.
\]
Integrating twice over $s\le x\le\rho$ gives
\begin{equation}\label{eq:large-prime-area}
 2\int_s^\rho\{(2x-1)_++(x-1)_+\}\,\dd x
 =3\rho^2-4\rho+\frac32=\frac{19}{16}.
\end{equation}
There is no factorial denominator prime beyond $2R-2$, and integrality
of the monomial matrix proves that no further denominator prime is needed.
Any additional numerator content can only decrease $Q_N$.

\subsection{Boundary primes and passage to the limit}
We spell out this passage because changing an endpoint convention at a
discontinuity can change a local exponent by order $N$.
It is not legitimate to call every such pointwise change $O(1)$.

Let $Y=N^{9/10}$. The small-prime estimate of
\cref{prop:small-primes} gives total cost $o(N^2)$ for $p\le Y$.
For larger primes, its hypotheses $p>5$ and $p^2>4R+20$ hold for all
sufficiently large $N$. Fix an integer $Q\ge40$ and restrict first to
$p/N\ge1/Q$. There are finitely many floor boundaries on this compact
interval. Each finite formula involves floors with numerators of the form
$cN+O(1)$. It can disagree with its limiting floor only when $p$ lies within
$O(1)$ of the corresponding fixed multiple of $N$. This leaves only $O(1)$
integer candidates per boundary.

For these candidates use the uniform entry bound
$v_p(F_N)\ge-6r$ from \cref{prop:local-kernel}. Their total cost is
$O(N\log N)=o(N^2)$. Elsewhere on the fixed interval the exponent is
$Ne(p/N)+O(1)$, with the profile specified above. Its finitely many
affine pieces cause no difficulty in prime summation.

We use the prime number theorem \cite[\S27.12]{dlmf} in the form
\begin{equation}\label{eq:pnt}
 \vartheta(x):=\sum_{p\le x}\log p=x+o(x).
\end{equation}
It implies, for every bounded piecewise continuous function $e$ of compact
support in $(0,\infty)$,
\[
 \frac1{N^2}\sum_p Ne(p/N)\log p\longrightarrow
 \int_0^\infty e(\beta)\,\dd\beta.
\]
For step functions this is exactly \eqref{eq:pnt} at finitely many fixed
endpoints; upper and lower step approximations prove the assertion for
the functions used here, including their signed values. The sum of the
$O(1)$ exponent errors is $O(N)$ on a fixed compact interval.

The primes $Y<p<N/Q$ have total cost at most
$6r\,\vartheta(N/Q)$ by the uniform bound. After division by $N^2$ the
limsup is at most $6\delta/Q$. Therefore we may first let $N\to\infty$
for fixed $Q$, then let $Q\to\infty$.
The profile \eqref{eq:rank-profile} is bounded between $0$ and $6\delta$,
so its omitted integral tends to zero as well. No prime number theorem
for intervals of shrinking relative length is being assumed.

\begin{proposition}\label{prop:arithmetic-rate}
The primitive normalizer satisfies
\begin{equation}\label{eq:arithmetic-rate}
 \limsup_{N\to\infty}\frac{\log Q_N}{N^2}\le C_{\arith},
\end{equation}
where
\begin{align}
 C_{\arith}
 &:=\frac{969}{6400}
   +\int_{1/40}^{9/20}e_{\mathrm{rank}}(\beta)\,\dd\beta
   +\frac{901}{600}+\frac{19}{16}\notag\\
 &=\frac{\scriptstyle
 191258319415140413215011151671222972968539453509572803845881093658106414835114844023}
 {\scriptstyle
 35352223433977492021442108242640841605890021236208954246644376534746108587827200000}
 <\frac{54101}{10000}.\label{eq:arithmetic-exact}
\end{align}
\end{proposition}
\begin{proof}
For every prime choose the integer exponent supplied by the applicable
finite local bound, using the boundary and small-prime choices just
described. The rational product of these prime powers clears every
coefficient of $F_N$. It is an integer multiple of $Q_N$, by the definition
of primitivity, so it is at least $Q_N$. Summing its logarithm and applying
the preceding limit argument proves the first line of
\eqref{eq:arithmetic-exact}.

For the remaining integral, the only initial breakpoints are
$1/q$, $2/(2q+1)$, and $2\rho/k$ in $[1/40,9/20]$.
Between consecutive such points, $q,\epsilon,D$ are constant. The only
additional breakpoint is the crossing
\[
 \beta=\frac{D+1-\delta}{(D+1)(q+\epsilon/2)}
\]
if it lies in that interval; it is where the two arguments of the minimum
in \eqref{eq:rank-profile} agree. These rules exhaust all changes of its
affine formula. Exact integration gives $228$ pieces on this interval.
Together with the seven displayed upper pieces there are $235$ pieces.
The code in \cref{app:certificate-code} evaluates their rational integrals
and verifies both the fraction and its strict comparison.
\end{proof}

% END SHARED CORE 11_arithmetic

% BEGIN SHARED CORE 12_certificate
\section{The exact rational comparison certificate}\label{sec:certificate}

Define
\[
 C_{\an}=-\int_0^\delta f_s(u)\,\dd u,\qquad g(u)=-f_s(u).
\]
We now give a lower bound for $C_{\an}$ using rational arithmetic only.
Every transcendental expression in this finite calculation is replaced by
a rational one-sided bound whose direction is proved here.

\begin{lemma}\label{lem:concavity}
The function $g$ is positive and continuous on $[0,\delta]$, and strictly
concave on $(0,\delta]$.
\end{lemma}
\begin{proof}
Differentiating \eqref{eq:ell-explicit}, or the equivalent expression in
the proof of \cref{prop:analytic-rate}, gives
\[
 f_s'(u)=4\log\frac{b}{4(s+u)},\qquad b=b_s(u).
\]
Write $v=s+u_s(b)$. Direct calculation gives
\[
 2(v-bu_s'(b))
 =1-\frac1{\sqrt{1+b^2}}+\frac{2s^2}{\sqrt{s^2+b^2}}>0.
\]
Since $u_s'(b)>0$,
\[
 f_s''(u)=\frac{4(v-bu_s'(b))}{b\,u_s'(b)\,v}>0
 \qquad(u>0).
\]
Also $v>b$: this follows immediately from
$\sqrt{1+b^2}<1+b$ and $\sqrt{s^2+b^2}>b$ in \eqref{eq:u-of-b}.
Hence $f_s'(u)<0$. Together with $f_s(0)=-4s\log2<0$, this proves
positivity of $g$ and the assertions about its shape.
\end{proof}

For a concave function the chord between any two graph points lies below
the graph. Thus
\begin{equation}\label{eq:trapezoid-lower}
 \int_{u_{j-1}}^{u_j}g(u)\,\dd u
 \ge\frac{u_j-u_{j-1}}2\{g(u_{j-1})+g(u_j)\}.
\end{equation}
Use the rational grid
\begin{equation}\label{eq:b-grid}
 b_j=\frac{96156889}{10^8}\frac{j}{2048},
 \quad 0\le j\le2048,\qquad u_j=u_s(b_j).
\end{equation}
The program verifies $u_{2048}<\delta$ by rational bounds. The omitted
terminal interval contributes positively and may be discarded.

\subsection{Square roots and logarithms with directed bounds}
Put $M=10^{30}$. For a nonnegative rational number $x$, let
$a=\operatorname{isqrt}(\floor{xM^2})$, where
$\operatorname{isqrt}(n)$ is the greatest integer whose square is at most
the nonnegative integer $n$. Then
\[
 a/M\le\sqrt x\le(a+1)/M.
\]
The code also checks both squared inequalities exactly.

For $1\le x\le2$ put $z=(x-1)/(x+1)$. Integration of the geometric series
for $2/(1-z^2)$, or differentiation followed by evaluation at $z=0$, gives
\[
 \log x=2\sum_{j=0}^\infty\frac{z^{2j+1}}{2j+1}.
\]
All terms are nonnegative. The first $24$ terms therefore give a lower
bound. The program rounds $z$, $z^2$, each successive odd power, and each
summand down to the $M^{-1}$ grid. Because every operation involves
nonnegative quantities, these rounds preserve a lower bound by induction.
For $x>2$, division by powers of two reduces it to $[1,2]$, and the same
lower bound for $\log2$ is added the appropriate number of times.
There is no need to bound the positive remainder from above.

For $b>0$, write $u=u_s(b)$. The expression used is
\begin{align}
 g(u)={}&4u\log\frac{4(s+u)}b
  +4s\log\frac{4(s+u)}{s+\sqrt{s^2+b^2}}
  +2\log\frac{1+\sqrt{1+b^2}}2.\label{eq:g-positive}
\end{align}
All three logarithms are nonnegative. Indeed $s+u>b$, and, if
$c=\sqrt{s^2+b^2}$, then
$s+u=(b+1-\sqrt{1+b^2})/2+c>c$ and $s\le c$.
The third assertion is immediate.
The prefactors and numerators can therefore be bounded downward, and
the second denominator upward, to obtain a lower bound.
The rational arguments supplied to the logarithm routine are checked
to be at least $1$.

At each $b_j$, let $u_j^-\le u_j\le u_j^+$ be the computed rational
interval and let $g_j^-\le g(u_j)$ be the positive rational lower bound.
For $j=0$ use $u_0^-=u_0^+=0$ and
$g_0^-=4s$ times the lower bound for $\log2$.
The program verifies $u_j^--u_{j-1}^+>0$ for every $j$.
It then sums
\begin{equation}\label{eq:rational-trapezoids}
 L=\sum_{j=1}^{2048}
   \frac{u_j^- -u_{j-1}^+}{2}(g_{j-1}^-+g_j^-).
\end{equation}
Every term is bounded above by its true trapezoid in
\eqref{eq:trapezoid-lower}; the omitted terminal integral is positive.

\begin{proposition}\label{prop:exact-gap}
The finite rational calculation proves
\[
 C_{\an}\ge L>5.4133705070>\frac{54133}{10000},
 \qquad
 C_{\an}-C_{\arith}>\frac{32}{10000}.
\]
\end{proposition}
\begin{proof}
The preceding inequalities prove $C_{\an}\ge L$. The complete
standard-library Python program in \cref{app:certificate-code} constructs
$L$ using only integers and rational numbers and checks the displayed
comparisons. Its arithmetic-profile computation is justified by the
exhaustive breakpoint argument in \cref{prop:arithmetic-rate}.
In particular, sampling is used only to select an affine branch whose
possible breakpoints have already been proved complete; a plot or a
numerical guess is not used to infer that completeness.
\end{proof}

The program is a small exact-arithmetic certificate. It is not a
proof-assistant development, and its successful execution does not
certify the analytic and local-arithmetic lemmas in the preceding
sections. These are separate mathematical arguments presented in full.

% END SHARED CORE 12_certificate

% BEGIN SHARED CORE 13_contradiction
\section{The positive-integer contradiction}\label{sec:contradiction}

\begin{proof}[Completion of the proof of \cref{thm:main}]
Suppose that $\zeta(5)=a/q_*$ with $a\in\Z$ and $q_*\in\Z_{>0}$.
By construction $Q_NF_N(X)\in\Z[X]$ and $\deg F_N\le N$.
Consequently
\[
 A_N:=q_*^N Q_NF_N(a/q_*)\in\Z.
\]
The Gram determinant is strictly positive at $\zeta(5)$, and both
normalizing factors are positive. Thus $A_N\ge1$ for every $N$ under
consideration.

On the other hand, \cref{prop:analytic-rate,prop:arithmetic-rate,prop:exact-gap}
give
\[
 \limsup_{N\to\infty}\frac{\log A_N}{N^2}
 \le C_{\arith}-C_{\an}<-\frac{32}{10000}.
\]
Here $N\log q_*/N^2\to0$, since the supposed denominator $q_*$ is fixed.
It follows that $A_N\to0$, contradicting $A_N\ge1$.
\end{proof}

\subsection{Scope of the conclusion and of its checks}
The logical inputs are the positive functional, the coefficientwise
local estimates, the explicit polynomial norm construction, the prime
number theorem, and the finite rational comparison proved in this
manuscript. No earlier lattice-covering proposal, stronger arithmetic
profile refinement, or external candidate proof of $\zeta(5)$ is an input.
The chosen window $s=1/20$, $\rho=5/4$ is sufficient; no optimization or
claim of best constants is needed.

This is a fully written \emph{candidate proof}, with internal mathematical
audits and a reproduced exact rational comparison. External peer review
has not been established. No Lean, Isabelle/HOL, Rocq/Coq, or other
proof-assistant certification of this particular argument is asserted.
The independently cited historical implementations establish precisely
the endpoints and trust boundaries listed in \cref{sec:history}.
The distinction matters even when two arguments have the same final
statement.

The present exposition also makes an endpoint qualification explicit:
floor-function jumps in \cref{sec:arithmetic} are treated by isolating
their finitely many nearby prime candidates before taking a limit.
This changes no displayed limiting constant, but it avoids an incorrect
pointwise interpretation of adjacent closed intervals in an asymptotic
profile table.

% END SHARED CORE 13_contradiction

% BEGIN JOURNAL ADDITION higher_odd
\section{Outlook for \texorpdfstring{$\zeta(7)$}{zeta(7)} and higher odd values}\label{sec:higher-odd}

The construction extends to higher odd arguments, but the strict
arithmetic--analytic comparison does not extend with the present window.
We make this distinction explicit. The available estimates for the primitive
coefficient normalizer are too costly to close the comparison; existence and
positivity of the determinant still hold. Related distinctions between rational approximation and the
inequalities needed for an irrationality conclusion occur in
\cite{prevostrivoal2021,zudilin2017}.

\subsection{The construction for a fixed higher odd argument}

Fix an integer $h\ge3$. All limits below keep $h$ fixed; they make no
assertion for an argument $2h+1$ that grows with $N$. Retain
$f(x)=(e^{2\pi x}-1)^{-1}$ and define a positive measure $\mu_{2h+1}$ by
\begin{equation}\label{eq:higher-measure}
 \int g(t)\,\dd\mu_{2h+1}(t)
 =\frac{2}{(2h)!}\int_0^\infty
       g(x^2)x^{2h+1}f^{(2h)}(x)\,\dd x.
\end{equation}
The formula initially applies to nonnegative measurable $g$.
The even derivative of the exponential series is positive. Its density
extends to $1/\pi$ at zero and is bounded by
$C_h(1+x)^{2h+1}e^{-2\pi x}$. Thus every polynomial moment exists.
The same calculation as in \cref{lem:positive-measure} gives
\begin{equation}\label{eq:higher-moments}
 \int t^j\,\dd\mu_{2h+1}(t)
 =\frac{(-1)^j}{(2h)!}
       \prod_{a=3}^{2h+1}(2j+a)\,\mathrm B_{2j+2}
 \in\Q,\qquad j\ge0.
\end{equation}

Write $\zeta(2h+1,y)=\sum_{n\ge0}(n+y)^{-(2h+1)}$ for $y>0$.
The Binet calculation of \cref{lem:stieltjes}, with $2h$ integrations
by parts, yields
\begin{equation}\label{eq:higher-stieltjes}
 \int\frac{\dd\mu_{2h+1}(t)}{t+y^2}
 =y^{2h}\zeta(2h+1,y)-\frac1{2y}-\frac1{2h}.
\end{equation}
Indeed, division of $x^{2h+1}$ by $x^2+y^2$ leaves
$(-1)^hy^{2h}x/(x^2+y^2)$; harmonicity contributes the same sign
when $2h$ derivatives in $x$ are replaced by derivatives in $y$.
The polynomial quotient disappears under these derivatives, and all
integration-by-parts boundary terms vanish as before.

Consequently the formal functional $\J_{2h+1,X}$ has the rational moments
in \eqref{eq:higher-moments} and the affine simple-fraction values
\begin{equation}\label{eq:higher-affine}
 \J_{2h+1,X}\!\left(\frac1{t+i^2}\right)
 =i^{2h}X-i^{2h}H_{i-1}^{(2h+1)}-\frac1{2i}-\frac1{2h},
 \quad H_n^{(2h+1)}=\sum_{a=1}^n a^{-(2h+1)}.
\end{equation}
Replace $\J_X$ by $\J_{2h+1,X}$ in \eqref{eq:window-pencil}, and denote
the resulting determinant and primitive normalizer by $F_N^{(2h+1)}$ and
$Q_N^{(2h+1)}$. Then
\[
 F_N^{(2h+1)}\in\Q[X],\qquad
 \deg F_N^{(2h+1)}\le N-m+1,\qquad
 F_N^{(2h+1)}(\zeta(2h+1))>0.
\]
Thus rationality $\zeta(2h+1)=a/b$, $b>0$, would again give the positive
integer $b^NQ_N^{(2h+1)}F_N^{(2h+1)}(a/b)$.

For fixed $h$, the extra polynomial factor in the density bound changes
only the lower-order errors in the norm estimates of
\cref{sec:discretization}. It does not change the external field $V_s$.
For the same fixed $s=m/N$ and $\rho=R/N$, the analytic rate therefore
remains the same explicitly defined $C_{\an}(s,\rho)$:
\begin{equation}\label{eq:higher-analytic-rate}
 \limsup_{N\to\infty}\frac1{N^2}
   \log F_N^{(2h+1)}(\zeta(2h+1))\le-C_{\an}(s,\rho).
\end{equation}
The constants in the error terms may depend on $h$.

\subsection{The changed local arithmetic}

Here are the precise replacements in the local estimates. It suffices
to use primes satisfying
\[
 p>2h+1,\qquad p^2>4R+4h+20.
\]
Enlarging this fixed good-prime cutoff has no effect on the limiting
prime sum. Write $N=qp+\sigma$, $\epsilon=\floor{2\sigma/p}$, and
$D=\floor{(2R-2)/p}$; let $w_p$ be the number of maximal reflected
classes in \eqref{eq:maximal-class-count}. The restriction-kernel proof
gives
\begin{equation}\label{eq:higher-rank-bound}
 v_p(F_N^{(2h+1)})\ge
 -(2h+\epsilon)r-\min\{r,(D+1)w_p\}.
\end{equation}
To see the changed exponent, the Bernoulli identity is now
\[
 \J_{2h+1,X}(q(t))=
 \frac1{(2h)!}\Bfun\!\left[
 \frac{\mathrm d^{\,2h-1}}{\mathrm d x^{2h-1}}
       \{x^{2h+1}q(-x^2)\}\right].
\]
There are $2h-1$ derivatives followed by one Bernoulli evaluation.
The residue harmonic sum costs at most $2h+1$ powers of $p$.
Both parts therefore admit the same restriction-kernel argument.
For the small primes the derivative count is still fixed, so their
total contribution remains $o(N^2)$.

There is a compensating improvement in the polynomial quotient.
In \eqref{eq:higher-moments}, the first $2h-1$ possible Bernoulli
poles cancel: at $2j+2=a(p-1)$ the product contains the factor $ap$,
for $1\le a\le2h-1$. The first uncanceled possible index is
$j=hp-h-1$. Thus the terminal fractional block of the quotient has size
\begin{equation}\label{eq:higher-terminal-block}
 t_{p,h}=\min\{r,\pos{2R-N-hp+h}\}.
\end{equation}

For a nonzero reflected class with $c$ high nodes, selecting $\ell$
atoms, of which $b$ are high, has cost at most
\[
 (2h+1)b+\ell(c+2-\ell).
\]
If $1\le c\le2h$, its maximum is $(2h+3)c$; an affected class has maximum
$(2h+1)c$ after the common-root gain. The deficit from either maximum
is at least one for each omitted atom below the maximizing count.
Classes with no high node are absorbed into the integral part, as in
\cref{lem:large-prime-integrality}.
These observations justify the replacements below, including the
shared dimension of the atoms and the complementary quotient minor.

\subsection{The original window does not give the required comparison}

Keep exactly
$s=1/20$, $\rho=5/4$, $\delta=6/5$, and the split at $\beta=9/20$
used in \cref{sec:arithmetic}. Let $C_{\arith}$ be the exact rational
number in \eqref{eq:arithmetic-exact}. Applying the preceding local
bounds with that same split gives a valid upper bound
\begin{equation}\label{eq:higher-arithmetic-rate}
 \limsup_{N\to\infty}\frac{\log Q_N^{(2h+1)}}{N^2}
 \le C_{\arith}^{(h)}=C_{\arith}+\Delta_h,
\end{equation}
where
\begin{equation}\label{eq:higher-rate-increment}
 \Delta_h=
 \frac{553}{400}(h-2)-\frac9{100}
 +\frac{\pos{3/2-9h/20}^{2}}{2h}.
\end{equation}
This is an exact rational expression, not a numerical fit.
On $0<\beta<9/20$, the rank profile increases by
$2(h-2)\delta$. On $9/20<\beta<1$, the atomic cost increases by
$2(h-2)(1-\beta)$, while the quotient contribution changes from
$(3/2-2\beta)_+$ to $(3/2-h\beta)_+$. The factorial charge and
the common-root gain do not change. Integrating these differences
gives \eqref{eq:higher-rate-increment}; the small-tail bound changes
by the integral of the same constant increment. For $p>N$ the
monomial matrix remains integral, so that part is unchanged as well.

\begin{center}
\begin{tabular}{@{}lll@{}}
\toprule
Target & Exact increment $\Delta_h$ & $C_{\arith}^{(h)}$ (rounded)\\
\midrule
$\zeta(7)$  & $1037/800$ & $6.7063289381$\\
$\zeta(9)$  & $107/40$   & $8.0850789381$\\
$\zeta(11)$ & $1623/400$ & $9.4675789381$\\
\bottomrule
\end{tabular}
\end{center}

The analytic rate in this table's window is unchanged. For orientation,
numerical integration of its explicit formula gives
$C_{\an}\approx5.4133707888$; accordingly the numerical difference
$C_{\an}-C_{\arith}^{(3)}$ is approximately $-1.2929581493$.
These decimal evaluations have no certificate role. The following
elementary upper bound proves that the displayed rates fail the
required comparison, independently of numerical quadrature.

\begin{proposition}\label{prop:higher-window-failure}
For this window, $C_{\an}<6<C_{\arith}^{(h)}$ for every $h\ge3$.
\end{proposition}
\begin{proof}
Retain $g(u)=-f_s(u)$ from \cref{sec:certificate}. It is increasing
and concave by \cref{lem:concavity} and its proof. Concavity, applied
to the pair $u,\delta-u$ and integrated, gives
\[
 C_{\an}\le\delta g(\delta/2).
\]
Direct squaring of rational bounds for the square roots in
\eqref{eq:u-of-b} gives
\[
 \frac35<u_s(1/2)<\frac{129}{200},\qquad
 \sqrt{1+(1/2)^2}<\frac98,\qquad
 \sqrt{s^2+(1/2)^2}>\frac12.
\]
Thus $s+u_s(1/2)<7/10$, and \eqref{eq:g-positive} implies
\begin{align*}
 g(u_s(1/2))
 &<\frac{129}{50}\log\frac{28}{5}
      +\frac15\log\frac{56}{11}+2\log\frac{17}{16}\\
 &<\frac{129}{50}\frac74+\frac15\frac53+\frac18
   =\frac{373}{75}<5.
\end{align*}
The first two logarithmic bounds follow respectively by retaining
terms through degree five in the positive series for $e^{7/4}$ and
through degree four for $e^{5/3}$; the third uses $\log(1+x)<x$.
Monotonicity of $g$ now gives $C_{\an}<\delta\cdot5=6$.
Finally, the exact rational value of $C_{\arith}$ and
\eqref{eq:higher-rate-increment} give
$C_{\arith}^{(3)}>67/10$, and the increments increase with $h\ge3$.
\end{proof}

This proposition concerns the explicitly derived rate bounds.
In particular, \eqref{eq:higher-arithmetic-rate} is an upper bound
for the primitive normalizer, not a lower bound. Its failure to beat
the analytic rate proves neither rationality of a higher odd value
nor impossibility of a sharper analysis of this determinant family.
No alternate window is asserted here to close the gap, and no
numerical optimization is used as evidence for such an assertion.

\subsection{The next structural questions}

\paragraph{First, use the available arithmetic bounds jointly.}
The split at $9/20$ was chosen for the present $\zeta(5)$ estimates;
it need not be retained for $\zeta(7)$. One can take the smaller of
the restriction-kernel and monomial bounds wherever both apply,
keeping the actual shared dimension and the signed factorial charge,
and then vary $s$ and $\rho$. There is also a justified larger
monomial range: for $\zeta(7)$, $q\le3$ gives at most six high nodes
per reflected class, so the preceding maximum and deficit estimates
apply for $p>N/4$, subject to $p>2m$. For $m\ge1$, the zero residue class has at
most three atoms, each of valuation at least $-1$; handling them
separately costs only $O(1)$ per such prime and hence $o(N^2)$ in
the prime sum. Extending this comparison is a precise first task,
not an assertion that a better window succeeds.

\paragraph{Second, change the common-root multiplicity.}
A concrete different subspace is spanned by
\begin{equation}\label{eq:higher-powered-window}
 B_m(t)^{a-1}B_{m+j}(t),\qquad 0\le j<r,\qquad a\ge2.
\end{equation}
Its monomial weight contains $C_m^{2a}/D_N$. For $p>2m$, each
surviving affected atom gains $2a$ powers of $p$, whereas the extra
factorial normalization is a $p$-unit. The largest polynomial degree,
however, increases to $R+(a-1)m-1$, and for $\zeta(7)$ the terminal
quotient block becomes
\[
 \min\{r,\pos{2R+2(a-1)m-N-3p+3}\}.
\]
The dimension accounting must also change. With $a=2$, an affected
class with $c$ high nodes has cost bounded by
$7\min(\ell,c)+\ell(c-2-\ell)$.
For $c=6$ its maximum $30$ occurs at both $\ell=5$ and $\ell=6$:
one omitted atom is free, and that extra available dimension must
be allowed in the complementary quotient minor.

There is a further useful local effect. For $p>2m$, the additional
factor $B_m$ reduces to zero on every affected residue class, and
to a nonzero constant on every unaffected class. Thus affected
maximal classes disappear from the leading restriction image;
on the remaining classes the original restriction-degree bound
still applies. These are genuine arithmetic gains for the new
subspace. Completing this route requires combining them with a
new analytic estimate for the changed external field
\[
 2\pi y+I_1(y)-2aI_s(y).
\]
The potential constructed for $a=1$ cannot simply be reused.
The effects of the larger quotient, the altered dimension count,
and the changed analytic field must all be included before a
numerical margin has mathematical significance.

Finally, rational rescaling alone cannot repair a negative margin.
For $c>0$ rational, primitivity gives
\[
 Q(cF)=Q(F)/c,\qquad Q(cF)\,cF=Q(F)F.
\]
A rational change of basis within the same polynomial subspace
multiplies its Gram determinant by a rational square and has the
same cancellation. Therefore improvements must sharpen the bounds
for this invariant quantity or change the underlying subspace or
rational-function family. The repeated-root construction above
does change the subspace; whether its gains suffice for $\zeta(7)$,
or for any higher individual odd value, remains an open question
within the present investigation.

% END JOURNAL ADDITION higher_odd

\appendix
\section{Complete exact comparison program}\label{app:certificate-code}
The following standalone program uses Python~3 and its standard library.
The reproduced run used CPython~3.12.14. Run it with
\texttt{python3 certificate\_explicit.py}; no optional package or network
access is needed. Its assertions check the finite comparison described in
Section~\ref{sec:certificate}. All displayed decimal thresholds in those
assertions are rational numbers.
\begin{lstlisting}
"""Exact rational certificate accompanying the Stage 60 manuscripts.

Python 3, standard library only. No floating-point arithmetic is used.
The program verifies the final inequalities conditional on the local and
analytic lemmas proved in the manuscripts. It is not a formalization of them.
"""
from fractions import Fraction as F
from math import isqrt

S = F(1, 20)
RHO = F(5, 4)
DELTA = RHO - S
LEFT = F(1, 40)
RIGHT = F(9, 20)
ZERO = F(0)
PRECISION = 10**30


def floor_data(beta):
    q = 1 // beta
    epsilon = int(2 * (1 - q * beta) >= beta)
    degree = (2 * RHO) // beta
    return q, epsilon, degree


def rank_profile(beta):
    q, epsilon, degree = floor_data(beta)
    width = 1 - (q + F(epsilon, 2)) * beta
    return (4 + epsilon) * DELTA + min(DELTA, (degree + 1) * width)


def exact_arithmetic():
    points = {LEFT, RIGHT}
    for q in range(1, 41):
        points.update((F(1, q), F(2, 2 * q + 1)))
    for k in range(1, 102):
        points.add(2 * RHO / k)
    points = sorted(x for x in points if LEFT <= x <= RIGHT)
    pieces = []
    for a, b in zip(points, points[1:]):
        q, epsilon, degree = floor_data((a + b) / 2)
        crossing = ((degree + 1 - DELTA)
                    / ((degree + 1) * (q + F(epsilon, 2))))
        cuts = [a, crossing, b] if a < crossing < b else [a, b]
        for left, right in zip(cuts, cuts[1:]):
            mid = (left + right) / 2
            width = 1 - (q + F(epsilon, 2)) * mid
            if (degree + 1) * width >= DELTA:
                intercept, slope = (5 + epsilon) * DELTA, ZERO
            else:
                intercept = (4 + epsilon) * DELTA + degree + 1
                slope = -(degree + 1) * (q + F(epsilon, 2))
            assert intercept + slope * mid == rank_profile(mid)
            area = (right - left) * (intercept + slope * mid)
            pieces.append((left, right, intercept, slope, area))
    assert pieces[0][0] == LEFT and pieces[-1][1] == RIGHT
    assert all(a[1] == b[0] for a, b in zip(pieces, pieces[1:]))
    assert len(pieces) == 228
    upper_rows = [
        (F(9, 20), F(1, 2), F(81, 5), -24),
        (F(1, 2), F(21, 40), F(139, 10), -15),
        (F(21, 40), F(5, 8), F(16), -19),
        (F(5, 8), F(3, 4), F(27, 2), -15),
        (F(3, 4), F(5, 6), F(12), -13),
        (F(5, 6), F(19, 20), F(19, 2), -10),
        (F(19, 20), F(1), F(38, 5), -8),
    ]
    upper_area = sum(((b - a) * (c + d * (a + b) / 2)
                      for a, b, c, d in upper_rows), ZERO)
    assert upper_area == F(901, 600)
    small_tail = (4 * DELTA + RHO) * LEFT + LEFT**2 / 4
    large_primes = 3 * RHO**2 - 4 * RHO + F(3, 2)
    assert small_tail == F(969, 6400)
    assert large_primes == F(19, 16)
    total = small_tail + large_primes + upper_area
    total += sum((row[4] for row in pieces), ZERO)
    expected = F(
        191258319415140413215011151671222972968539453509572803845881093658106414835114844023,
        35352223433977492021442108242640841605890021236208954246644376534746108587827200000)
    assert total == expected
    assert total < F(54101, 10000)
    return total, len(pieces) + len(upper_rows)


def sqrt_bounds(x):
    """Outward rational square-root bounds of width 10^-30."""
    root = isqrt(x.numerator * PRECISION**2 // x.denominator)
    lower, upper = F(root, PRECISION), F(root + 1, PRECISION)
    assert lower * lower <= x <= upper * upper
    return lower, upper


def log_unit_lower(x, terms=24):
    """Positive atanh series with every operation rounded down."""
    assert 1 <= x <= 2
    z = (x - 1) / (x + 1)
    iz = z.numerator * PRECISION // z.denominator
    iz2 = iz * iz // PRECISION
    power, total = iz, 0
    for j in range(terms):
        total += 2 * power // (2 * j + 1)
        power = power * iz2 // PRECISION
    return F(total, PRECISION)


LOG_TWO_LOWER = log_unit_lower(F(2))


def log_lower(x):
    assert x >= 1
    exponent = 0
    while x > 2:
        x /= 2
        exponent += 1
    return log_unit_lower(x) + exponent * LOG_TWO_LOWER


def u_and_g_lower(b):
    if b == 0:
        return ZERO, ZERO, 4 * S * LOG_TWO_LOWER
    a_lower, a_upper = sqrt_bounds(1 + b * b)
    c_lower, c_upper = sqrt_bounds(S**2 + b * b)
    u_lower = (b + 1 - 2 * S - a_upper + 2 * c_lower) / 2
    u_upper = (b + 1 - 2 * S - a_lower + 2 * c_upper) / 2
    assert 0 < u_lower <= u_upper
    v_lower = S + u_lower
    g_lower = (4 * u_lower * log_lower(4 * v_lower / b)
               + 4 * S * log_lower(4 * v_lower / (S + c_upper))
               + 2 * log_lower((1 + a_lower) / 2))
    assert g_lower > 0
    return u_lower, u_upper, g_lower


def exact_analytic_lower():
    endpoint = F(96156889, 10**8)
    subdivisions = 2048
    _, previous_upper, previous_g = u_and_g_lower(ZERO)
    total = ZERO
    for j in range(1, subdivisions + 1):
        b = endpoint * F(j, subdivisions)
        u_lower, u_upper, g_lower = u_and_g_lower(b)
        width_lower = u_lower - previous_upper
        assert width_lower > 0
        total += width_lower * (g_lower + previous_g) / 2
        previous_upper, previous_g = u_upper, g_lower
    assert previous_upper < DELTA
    assert total > F(54133705070, 10**10)
    assert total > F(54133, 10000)
    return total


def main():
    arithmetic, pieces = exact_arithmetic()
    analytic_lower = exact_analytic_lower()
    assert analytic_lower - arithmetic > F(32, 10000)
    print("EXACT COMPARISON CERTIFICATE PASS")
    print("Arithmetic affine pieces:", pieces)
    print("Arithmetic exact:", arithmetic)
    print("Arithmetic < 54101/10000 = 5.4101")
    print("Analytic lower bound > 54133/10000 = 5.4133")
    print("Certified strict margin > 32/10000 = 0.0032")
    print("Analytic lower bound, rounded down at 10 decimal places:",
          analytic_lower.numerator * 10**10 // analytic_lower.denominator,
          "/ 10000000000")
    print("This checks the comparison, not the underlying mathematical lemmas.")


if __name__ == "__main__":
    main()

\end{lstlisting}

\begingroup
\raggedright
\begin{thebibliography}{99}
% Bibliography items for the shared historical introduction.
% Insert inside the main document's thebibliography environment.
% Research cutoff: 13 November 2026. No formal development was rebuilt here.

\bibitem{euler1740}
L. Euler, \emph{De summis serierum reciprocarum}, Commentarii academiae
scientiarum Petropolitanae \textbf{7} (1740), 123--134, Enestr\"om E41;
written 1735.  Euler Archive,
\url{https://scholarlycommons.pacific.edu/euler-works/41/}.

\bibitem{lindemann1882}
F. Lindemann, \emph{\"Uber die Zahl $\pi$}, Mathematische Annalen
\textbf{20} (1882), 213--225.

\bibitem{dlmf}
NIST Digital Library of Mathematical Functions, \S25.6, \emph{Integer
Arguments}, especially formula 25.6.2;
\S5.9, \emph{Integral Representations}, especially formula 5.9.15;
and \S27.12, \emph{Asymptotic Formulas: Primes}.
\url{https://dlmf.nist.gov/25.6};
\url{https://dlmf.nist.gov/5.9};
\url{https://dlmf.nist.gov/27.12}; consulted 13 November 2026.

\bibitem{apery1979}
R. Ap\'ery, \emph{Irrationalit\'e de $\zeta 2$ et $\zeta 3$}, in
Journ\'ees Arithm\'etiques de Luminy, Ast\'erisque \textbf{61} (1979),
11--13. \url{https://www.numdam.org/item/AST_1979__61__11_0/}.

\bibitem{beukers1979}
F. Beukers, \emph{A note on the irrationality of $\zeta(2)$ and
$\zeta(3)$}, Bulletin of the London Mathematical Society \textbf{11}
(1979), no.~3, 268--272. \url{https://doi.org/10.1112/blms/11.3.268}.

\bibitem{rivoal2000}
T. Rivoal, \emph{La fonction z\^eta de Riemann prend une infinit\'e de
valeurs irrationnelles aux entiers impairs}, Comptes Rendus de
l'Acad\'emie des Sciences, S\'erie I, \textbf{331} (2000), 267--270.
\url{https://doi.org/10.1016/S0764-4442(00)01624-4};
\url{https://arxiv.org/abs/math/0008051v1}.

\bibitem{ballrivoal2001}
K. Ball and T. Rivoal, \emph{Irrationalit\'e d'une infinit\'e de valeurs
de la fonction z\^eta aux entiers impairs}, Inventiones Mathematicae
\textbf{146} (2001), no.~1, 193--207.

\bibitem{rivoal2002}
T. Rivoal, \emph{Irrationalit\'e d'au moins un des neuf nombres
$\zeta(5),\zeta(7),\ldots,\zeta(21)$}, Acta Arithmetica \textbf{103}
(2002), no.~2, 157--167; preprint 2001.
\url{https://arxiv.org/abs/math/0104221}.

\bibitem{zudilin2001}
W. Zudilin, \emph{One of the numbers $\zeta(5)$, $\zeta(7)$,
$\zeta(9)$, $\zeta(11)$ is irrational}, Russian Mathematical Surveys
\textbf{56} (2001), no.~4, 774--776.
\url{https://doi.org/10.1070/RM2001v056n04ABEH000427};
\url{https://www.mathnet.ru/eng/rm427}.

\bibitem{zudilin2001windows}
W. Zudilin, \emph{One of the eight numbers
$\zeta(5),\zeta(7),\ldots,\zeta(17),\zeta(19)$ is irrational},
Mathematical Notes \textbf{70} (2001), no.~3, 426--431; Theorems~2--3
give the shifted windows ending at 35 and 51.
\url{https://doi.org/10.1023/A:1012312315852};
\url{https://www.mathnet.ru/eng/mzm759}.

\bibitem{zudilin2001note}
W. Zudilin, \emph{Irrationality of values of zeta-function},
preprint, 26 April 2001. \url{https://arxiv.org/abs/math/0104249v1}.

\bibitem{fsz2019}
S. Fischler, J. Sprang, and W. Zudilin, \emph{Many odd zeta values are
irrational}, Compositio Mathematica \textbf{155} (2019), no.~5, 938--952.
\url{https://doi.org/10.1112/S0010437X1900722X};
\url{https://arxiv.org/abs/1803.08905}.

\bibitem{rivoalzudilin2020}
T. Rivoal and W. Zudilin, \emph{A note on odd zeta values}, S\'eminaire
Lotharingien de Combinatoire \textbf{81} (2020), article B81b, 13 pp.;
preprint 8 March 2018. \url{https://arxiv.org/abs/1803.03160v1};
\url{https://www.mat.univie.ac.at/~slc/wpapers/s81rivzud.html}.

\bibitem{laiyu2020}
L. Lai and P. Yu, \emph{A note on the number of irrational odd zeta
values}, Compositio Mathematica \textbf{156} (2020), 1699--1717.
\url{https://doi.org/10.1112/S0010437X20007307};
\url{https://arxiv.org/abs/1911.08458v2}.

\bibitem{laizhou2022}
L. Lai and L. Zhou, \emph{At least two of
$\zeta(5),\zeta(7),\ldots,\zeta(35)$ are irrational},
Publicationes Mathematicae Debrecen \textbf{101} (2022), nos.~3--4,
353--372. \url{https://doi.org/10.5486/PMD.2022.9252};
\url{https://arxiv.org/abs/2103.00904v2}.

\bibitem{fischler2026}
S. Fischler, \emph{Linear independence of odd zeta values using
Siegel's lemma}, Journal of the London Mathematical Society
\textbf{113} (2026), no.~4, e70535; published 20 April 2026.
\url{https://doi.org/10.1112/jlms.70535}.

\bibitem{lai2026count}
L. Lai, \emph{A note on the number of irrational odd zeta values, II}, Journal
de Th\'eorie des Nombres de Bordeaux \textbf{38} (2026), no.~2,
469--487; published 16 September 2026.
\url{https://doi.org/10.5802/jtnb.1368};
\url{https://arxiv.org/abs/2501.05321v2}.

\bibitem{lai2026small}
L. Lai, \emph{Small improvements to the Ball--Rivoal theorem and its
$p$-adic variant}, Frontiers in Combinatorics and Number Theory
\textbf{4} (2026), 1--44; published 8 October 2026.
\url{https://doi.org/10.3934/fcnt.2026024}.

\bibitem{lai2026linear}
L. Lai, \emph{The linear independence of $1$, $\zeta(3)$, and
$\zeta(4)$}, preprint, submitted 5 October 2026, revised 15 October
2026. \url{https://arxiv.org/abs/2610.06838v2}.

\bibitem{fauzan2026}
A. Fauzan, \emph{$\zeta(5)$ is irrational}, Zenodo preprint, version~1,
17 September 2026. \url{https://doi.org/10.5281/zenodo.22826419};
\url{https://zenodo.org/records/22826419}.

\bibitem{delsolar2026}
C. Del Solar, \emph{Zeta5: a Lean formalization of Fauzan's proof},
repository snapshot 23 September 2026, Lean 4.35.0-rc2.
\url{https://github.com/domino14/zeta5/tree/3da8eb1b91f9a90a1c0b7959e845c060dbd16256}.
Mathlib commit \texttt{065356127b1dc0016f66b7283ce0ce2c4055aa55}.
See the README, \texttt{Zeta5/lean-toolchain}, and
\texttt{Zeta5/Claim.lean}; reported endpoint axiom audit, not replayed
for this survey.

\bibitem{firsching2026}
M. Firsching, \emph{Zeta5}, repository snapshot 27 September 2026,
Lean 4.34.0-rc1.
\url{https://github.com/mo271/Zeta5/tree/7fe736760f4b96bfdb4334b68e3b3124ecbe10b0}.
Mathlib commit \texttt{f385f9fd18abed332ebcd3df239cb357e264aceb};
PrimeNumberTheoremAnd commit
\texttt{04d81b2a68ea3548b9990813fdc463c78511f412}.
Build record: \url{https://reservoir.lean-lang.org/@mo271/Apery}.

\bibitem{romik2026}
D. Romik, \emph{A Lean formalization of the irrationality of
$\zeta(5)$}, repository snapshot 25 September 2026, Lean 4.34.0.
\url{https://github.com/danromik/zeta5-irrationality/tree/fc2fefa77076e22e7e80b4c101ded5a98d53e14d}.
Mathlib commit \texttt{5ed2965256430c3649e86755f9576b54eca72435}.
The report \texttt{docs/CERTIFICATION.md} records checks of parent
commit \texttt{a87234f50f4c0d257a2aaed0894e6587c022e5ba}; the explicit
prime-number-theorem axiom remains part of the endpoint assumptions.

\bibitem{hyman2026}
A. R. Hyman, \emph{The irrationality of $\zeta(5)$}, Archive of
Formal Proofs, 26 October 2026, Isabelle/HOL development.
\url{https://isa-afp.org/entries/Zeta_5_Irrational.html}.
Endpoint and oracle report:
\url{https://isa-afp.org/browser_info/current/AFP/Zeta_5_Irrational/Zeta_5_Irrational.html}.
Exact Isabelle release not established from the inspected record.

\bibitem{devincenzis2026}
G. M. Devincenzis, \emph{At least one of
$\zeta(7),\zeta(9),\ldots,\zeta(21)$ is irrational}, Lean development,
September 2026; Lean 4.35.0-rc2.
\url{https://github.com/gmDevi/zeta-7-21-lean/tree/4e48690ab5f9d04cbf3cede2b657f8db91ef6878}.
Mathlib commit \texttt{065356127b1dc0016f66b7283ce0ce2c4055aa55}.
The repository reports registration PALOMAR-2026-09-30-000007, version~1,
for snapshot \texttt{a3b69fc43aea3dd140540f6e8380c7d0fecce028}; the
registry page itself was unavailable in this survey.

\bibitem{lin2026}
Y. Lin, \emph{Odd zeta irrationality}, repository snapshot
29 September 2026, Lean 4.35.0-rc3.
\url{https://github.com/CoolRmal/odd-zeta-irrationality/tree/4a9a71ccbd1d7ef229bf149e56f3f8d6c53416eb}.
Mathlib commit \texttt{3cb72cfd416d1b5ec4b930d67648ef036a5df24f}.
The inspected successful checker workflow instead fixes commit
\texttt{b1f6999fbcb0f7c6bdef92ed172e845d39f99732}:
\url{https://github.com/CoolRmal/odd-zeta-irrationality/actions/runs/36503793020}.

\bibitem{tang2026}
Q. Tang, \emph{Linear independence of $1$, $\zeta(2)$ and $\zeta(3)$
over $\Q$}, repository snapshot 30 September 2026, Lean 4.35.0-rc3.
\url{https://github.com/dtq1997/zeta2-zeta3-linear-independence/tree/669c92bf7a3728e0de9ad80c373bc29fb5ca3d58}.
Mathlib commit \texttt{c55e6e786f49471c72fbddbec5415808896aec1e}.
Palomar registration PALOMAR-2026-09-30-000030, version~1:
\url{https://palomar-registry.org/entry?id=PALOMAR-2026-09-30-000030&version=1}.
The indexed registry record identifies Lean toolchain commit
\texttt{470d5ce1400764999581fd26d5d72b00d990b0f4} and successful
Lean, NanoDa, and \texttt{con-ron} checks on 30 September 2026;
their proof-checking run was not replayed here.

\bibitem{jingweiliu2026}
J. Liu (Jingwei Liu), \emph{At least one of
$\zeta(7),\zeta(9),\ldots,\zeta(19)$ is irrational: a growing-determinant
proof with Lean certification}, version~1.0.0, 6 October 2026,
Lean 4.34.0-rc2. \url{https://doi.org/10.5281/zenodo.23191331}.
\url{https://github.com/AlexxxxxLiu/odd-zeta-seven-lean/tree/7a7e3a2895c468306d28ba584558aed6decf45aa}.
Mathlib commit \texttt{dc4b8d60d5edb3c493c3662126b1b7ccae7d67cf}.

\bibitem{krisman2026}
S. Krisman, \emph{One of $\zeta(7),\zeta(9),\ldots,\zeta(21)$ is
irrational}, preprint, \url{https://hexagonmath.org/2610.00516v2};
Lean 4.35.0 formalization reported complete 22 October 2026.
Inspected source snapshot, 11 November 2026:
\url{https://github.com/telaviv/zeta-irrationality/tree/d6ec102d61ed954169f7a8778e28f21a80bf6851}.
Mathlib commit \texttt{56a005855ea17e0b0958f73ec801ae140b0a1679}.

\bibitem{davidson2026}
M. Davidson, \emph{Four odd zeta values: a determinant argument},
Round~29 report, 15 October 2026; source snapshot 16 October 2026.
\url{https://github.com/morley988/four-odd-zeta/tree/fd23dd59e5d638c9684eb2d7776fb51e8b427b55}.
Lean 4.33.1; Mathlib commit
\texttt{0df444a360eaa60ab8c11dca51a86af692955474};
PrimeNumberTheoremAnd commit
\texttt{650d31264be65f4cd6e70c45d8b25d86d482a761}.
See \texttt{docs/VERIFICATION.md} and
\texttt{proof/Four\_Value\_Final\_Report\_2026-10-15.txt} for the
certificate-conditional endpoint and the eight numerical admissions.

\bibitem{chyzak2014}
F. Chyzak, A. Mahboubi, T. Sibut-Pinote, and E. Tassi,
\emph{A computer-algebra-based formal proof of the irrationality of
$\zeta(3)$}, Interactive Theorem Proving, Lecture Notes in Computer
Science \textbf{8558} (2014), 160--176.
\url{https://doi.org/10.1007/978-3-319-08970-6_11}.

\bibitem{mahboubi2021}
A. Mahboubi and T. Sibut-Pinote, \emph{A formal proof of the
irrationality of $\zeta(3)$}, Logical Methods in Computer Science
\textbf{17} (2021), no.~1, article~16, 16:1--16:25.
\url{https://doi.org/10.23638/LMCS-17(1:16)2021};
\url{https://arxiv.org/abs/1912.06611v6}.
The completed paper cites Coq 8.11.0.

\bibitem{aperysource}
Rocq Community, \emph{apery}, maintained source snapshot
25 April 2026.
\url{https://github.com/rocq-community/apery/tree/97f0c01bec4282c4504c0225a71ff0e75aeab115}.
The current source specifies Coq 8.20+; CI configurations for Coq 8.20
and Rocq 9.0/9.1 are version declarations, not fresh build results
reported by this survey.

\bibitem{eberl2019}
M. Eberl, \emph{The irrationality of $\zeta(3)$}, Archive of Formal
Proofs, 27 December 2019.
\url{https://isa-afp.org/entries/Zeta_3_Irrational.html}.
An inspected development-build receipt, dated 13 November 2026, names
Isabelle revision \texttt{559e5864efb7} and AFP revision
\texttt{62317ad25318}:
\url{https://devel.isa-afp.org/entries/Zeta_3_Irrational.html}.
The original 2019 Isabelle release was not established here.

\bibitem{liuzhangzhi2025}
J. Liu (Junqi Liu), J. Zhang, and L. Zhi, \emph{A formal proof of the
irrationality of $\zeta(3)$ in Lean 4}, preprint, revised 8 August 2025.
\url{https://arxiv.org/abs/2503.07625v2}.
Source snapshot 20 March 2026, Lean 4.18.0:
\url{https://github.com/ahhwuhu/zeta_3_irrational/tree/e8785315a01c8fbcddaa0fc03b3c8b29a61bc1f1}.
Mathlib commit \texttt{aa936c36e8484abd300577139faf8e945850831a}.

\bibitem{anand2026}
P. Anand, \emph{Zeta7 is Irrational}, announced Zenodo preprint,
23 September 2026, version~1,
\url{https://doi.org/10.5281/zenodo.22920911}.
The source claim is cited in the Devincenzis repository; the Zenodo
record itself was unavailable during this survey.  This entry is not
an endorsement of the argument.

\bibitem{anandobjection}
Public discussion of the arithmetic in Anand's claimed $\zeta(7)$
argument, issue~1 of the Devincenzis repository,
\url{https://github.com/gmDevi/zeta-7-21-lean/issues/1}; consulted
13 November 2026.

\bibitem{zudilin2017}
W. Zudilin, \emph{A determinantal approach to irrationality},
Constructive Approximation \textbf{45} (2017), no.~2, 301--310.
\url{https://doi.org/10.1007/s00365-016-9333-7};
\url{https://arxiv.org/abs/1507.05697v2}.

\bibitem{hayneszudilin2015}
A. Haynes and W. Zudilin, \emph{Hankel determinants of zeta values},
SIGMA \textbf{11} (2015), article~101, 5 pp.
\url{https://doi.org/10.3842/SIGMA.2015.101};
\url{https://arxiv.org/abs/1510.01901v3}.

\bibitem{prevost1996}
M. Pr\'evost, \emph{A new proof of the irrationality of $\zeta(2)$ and
$\zeta(3)$ using Pad\'e approximants}, Journal of Computational and Applied
Mathematics \textbf{67} (1996), no.~2, 219--235.
\url{https://doi.org/10.1016/0377-0427(95)00019-4}.

\bibitem{prevost2019}
M. Pr\'evost, \emph{Remainder Pad\'e approximants for the Hurwitz zeta
function}, Results in Mathematics \textbf{74} (2019), no.~1, article~51,
22~pp. \url{https://doi.org/10.1007/s00025-019-0976-0}.
Preprint: \url{https://arxiv.org/abs/1709.05389}.
The section and equation references in the text refer to the arXiv version.

\bibitem{prevostrivoal2021}
M. Pr\'evost and T. Rivoal, \emph{Diagonal convergence of the remainder
Pad\'e approximants for the Hurwitz zeta function}, Journal of Number
Theory \textbf{222} (2021), 346--361.
\url{https://doi.org/10.1016/j.jnt.2020.10.019};
\url{https://hal.science/hal-02498437}.

\bibitem{calegari2026}
F. Calegari, \emph{zeta(5) is irrational}, Persiflage,
24 September 2026, online mathematical exposition.
\url{https://galoisrepresentations.org/2026/09/24/zeta5-is-irrational/};
consulted 13 November 2026.
This reference is used for the comparison of constructions, not as an
input to the arithmetic or analytic estimates proved here.

\end{thebibliography}
\endgroup
\end{document}
